\documentclass[12pt,a4paper,fleqn,spanish,final,openright]{extbook}
\usepackage{amsbsy}
\usepackage{amssymb}
\usepackage{fontspec}
\setcounter{secnumdepth}{3}
\setcounter{tocdepth}{3}
\synctex=1
\usepackage[unicode=true,
 bookmarks=true,bookmarksnumbered=true,bookmarksopen=true,bookmarksopenlevel=4,
 breaklinks=false,pdfborder={0 0 1},backref=section,colorlinks=true]
 {hyperref}
\hypersetup{pdftitle={Clasificación de los grupos hasta el orden 22},
 pdfauthor={Julián Calderón Almendros},
 pdfsubject={Trabajo Fin de Grado del Grado de Matemáticas UNED},
 pdfkeywords={Álgebra, Grupos Finitos, Teoría de Números, Matemáticas Discretas}}

\makeatletter

\usepackage[spanish]{babel}
\usepackage{amsthm}
\usepackage{amsfonts}
\usepackage{amsxtra}
\usepackage{amscd}
\usepackage{bbm}
\usepackage{lplfitch}
\usepackage{abraces}
\usepackage{color}
\usepackage{dsfont}
\usepackage{mathrsfs}
\usepackage{cmll}
\usepackage{latexsym}
\usepackage[mathscr]{euscript}
\usepackage{color}
\usepackage{units}
\usepackage{wasysym}
\usepackage{nomencl}
% the following is useful when we have the old nomencl.sty package
\providecommand{\printnomenclature}{\printglossary}
\providecommand{\makenomenclature}{\makeglossary}
\makenomenclature
\usepackage{microtype}
\makeatletter
\newtheorem{thm}{Teorema}[chapter]
\newtheorem{example}[thm]{Ejemplo}
\newtheorem{exercise}[thm]{Ejercicio}
\newtheorem{axiom}[thm]{Axioma}
\newtheorem{cor}[thm]{Corolario}
\newtheorem{lem}[thm]{Lema}
\newtheorem{prop}[thm]{Proposici\'on}
\theoremstyle{definition}
\newtheorem{defn}[thm]{Definici\'on}
\theoremstyle{remark}
\newtheorem{rem}[thm]{Nota}
\newtheorem{fixed}[thm]{Fijaci\'on de Notaci\'on}
\newtheorem{formalproof}[thm]{Demostraci\'on Formal}
\newtheorem{solution}[thm]{Soluci\'on}

\makeatother

\usepackage{listings}
\usepackage{tikz}
%\usepackage{polyglossia}
%\setdefaultlanguage{spanish}
%\def\eqdeclaration#1{, véase la ecuación\nobreakspace(#1)}
%\def\pagedeclaration#1{, página\nobreakspace#1}
%\def\nomname{Nomenclatura}
%\renewcommand{\lstlistlistingname}{Listados de código}

\begin{document}
\tableofcontents{}

\listoffigures

\listoftables
%\lstlistoflistings{}\settowidth{\nomlabelwidth}{${\bold{G}}:={\langle{G,+,e}\rangle}$}
\printnomenclature{}

\title{Clasificación de grupos finitos no isomorfos hasta el orden 22}

\title{{\normalsize{}T.F.G. para el grado en Matemáticas de la UNED}}

\author{Julián Calderón Almendros}

\author{\textsc{\small{}TUTORIZADO POR} \textbf{Emilio Bujalance
García}}

\date{24/04/2018}
\maketitle

\part{Preliminares teóricos}

En estos capítulos que ahora escribo quiero dejar sentado los «hechos»
que tomaré como ya probados, mayormente provenientes de la Unidad
Didáctica I de la asignatura de Álgebra de 2º curso de la antigua
licenciatura en Matemáticas, tomo dedicado precisamente a la teoría
de grupos.\nomenclature{$\bold{G}:={\langle{G,*,u}\rangle}$}{Forma 
habitual de nombrar un grupo: el conjunto más la operación binaria, 
más la nularia del elemento neutro. Habitualmente como la operación es 
no conmuitativa, llamo al elemento neutro $u$ de unidad. El nombre del 
conjunto subyacente al grupo se nombra como todos los conjuntos de este 
mismo nivel y subconjuntos del mismo, en mayúsculas simples. Para nombrar 
el grupo propiamente dicho, nombro el conjunto subyacente con letras negritas.}
\nomenclature{${\bold{G}}:={\langle{G,+,e}\rangle}$}{Forma habitual de nombrar 
un grupo abeliano: el conjunto más la operación binaria, más la nularia del 
elemento neutro.Esta vez la operación es conmutativa,así que utilizamos notación 
aditiva y  llamo al elemento neutro $e$ de neutro.}

\chapter{Algunas definiciones}


Comenzaremos directamente. Los grupos cíclicos de un determinado orden
$n$ lo denominaremos $\mathbf{C_{n}}$. El grupo trivial será denotado
por $\mathbf{C_1}$. Los homomorfismos entre dos grupos, digamos
de $G$ en $L$serán denotados como $Hom(G,L)$. Si estos son biyecciones
se denotará como $Isom(G,L)$. Si los homomorfismos son de un grupo
en sí mismo pondremos $End(G)$ y si son biyecciones $Aut(G)$. También
serán de interés las biyecciones de un conjunto en sí mismo $Biy(S)$
que es además un grupo bajo la composición de funciones como operación
de grupo. Las operaciones entre grupos que utilizaremos serán $\rtimes$
para el producto semidirecto, dónde, si escribimos $N\rtimes H$ siendo
$N \trianglelefteq G$ subgrupo normal del grupo en el que están y
$H \preccurlyeq G$ simplemente subgrupo. Esta forma de nombrar el
producto directo es imprecisa, así que siendo $\varphi \in Hom(H,Aut(N))$,
el simbolismo completo será $N\rtimes_{\varphi }H$. Cuando ambos subgrupos
sean normales simplemente escribiremos $N\times  H$, que es el producto
directo. A su vez el producto semidirecto será considerado como un
caso particular del producto SZ, que solo requiere que ambos subconjuntos
sean subgrupos, así que necesitaremos dos homomorfismos, $\alpha \in Hom(N,Biy(H))$
y $\beta \in Hom(H,Biy(N))$, denotándose como $N_{\alpha \!}\bowtie_{\beta }H$
. Más adelante desarrollaremos las relaciones entre estos productos
y algunas de sus propiedades, y veremos las diferencias entre las
versiones externa e interna de ellos. Además, cuando un subgrupo de
$G$, pongamos $H \preccurlyeq G$, podemos definir dos relaciones
de equivalencia $\mathcal{R}_{H}$ y $\mathcal{R}^{H}$, y el conjunto
de particiones para una de ellas será $\nicefrac{G}{{\mathcal{R}}_{\mathrm{H}}}
\triangleq\lbrace gH \vert g\in G\rbrace$,
y para la otra 
$\nicefrac{G}{{\mathcal{R}}^{\mathrm{H}}}\triangleq\lbrace Hg|g\in G\rbrace$
. Cuando $H\trianglelefteq G$ lo denotaremos por $\nicefrac{G}{H}$
y además constituirá un grupo.

Cuando un grupo sea hamiltoniano usaremos el símbolo $\boldsymbol{Q}_{8}$ ya
que tiene alguna copia de él. Son importantes estos grupos porque
no se pueden construir a través de los productos antes citados (al menos en 
sus formas más elementales).

Hay aún dos productos importantes que citar: uno es el producto wreath,
denotado $N\wr H$, y que es un producto semidirecto a partir de
$\sigma \in Hom(H,Biy(N))$, dónde, si $h \in H$ tenemos que $\sigma_{h} \in Biy(N)$
que se forma de una manera especial. El otro producto que nos queda
es, sean $H$ y $K$ grupos, tales que tienen un subgrupo cada uno
de ellos $I \preccurlyeq \mathsf{Z}(H)$ y $J \preccurlyeq \mathsf{Z}(K)$
tales que podemos encontrar un isomorfismo $\delta \in Isom(I,J)$
con el que podemos crear $L\triangleq\lbrace(i,j) \vert i\in I,j\in J,\delta (i)*j=1\rbrace$.
Ahora el producto central es $\nicefrac{H\times K}{L}$.

De los subgrupos de un grupo dado nos interesan el ya mencionado $\mathsf{Z}(G)$
o centro, y en general $\mathsf{C}_{G}(H)$, siendo $H \preccurlyeq G$,
el normalizador del anterior subgrupo $\mathsf{N}_{G}(H)$. Para definir
los siguientes subgrupos, vamos a definir las primicipales familias
de subgrupos. 
$\mathcal{S}_{G}\triangleq\lbrace H \subseteq G\vert H \preccurlyeq G\rbrace$,
$\mathcal{N}_{G}\triangleq\lbrace H \in \mathcal{S}_{G}|H\trianglelefteq G\rbrace$. 
Desde aquí podemos definir 
$\mathcal{\overline N}_{G}(K)\triangleq\lbrace H \in \mathcal{N}_{G}\vert H\subseteq G\rbrace$, 
y la definición de un grupo normal 
${\overline{\mathsf{K}}}_{G}(K)\triangleq\bigcap\mathcal{\overline N}_{G}(K)$ 
y de otro 
${\underline{\mathsf{K}}}_{G}(K)\triangleq\left\langle\bigcup\mathcal{\underline N}_{G}(K)\right\rangle$.
También nos quedan los subgrupos característicos, invariantes ante
isomorfismos cualquiera, $ H \overset{\mathsf{{char}}}{\_} G $. Así definimos
$\mathcal{C}_{G}\triangleq\lbrace H \in \mathcal{N}_{G} \vert H \overset{\mathsf{{char}}}{\_} G$.
Además, será interesante el conjunto de los subgrupos maximales de $G$, 
$\mathcal{M}_{G} \triangleq
\lbrace H \in \left({\mathcal{S}_{G}}\setminus \lbrace G\rbrace\right) \vert
\exists K\in \left({\mathcal{S}_{G}}\setminus \lbrace G\rbrace\right)\; 
H \subseteq K \Rightarrow H=K\rbrace$. 
Para subgrupos minimales tenemos dualmente 
$\mathcal{P}_{G} \triangleq
\lbrace H \in \left({\mathcal{S}_{G}}\setminus \lbrace \mathsf{C}_{1}\rbrace\right) 
\vert \exists K \in \left({\mathcal{S}_{G}}\setminus \lbrace 
\mathsf{C}_{1}\rbrace\right)\; K \subseteq H \Rightarrow H=K\rbrace$. 
Llamamos subgrupo de Frattini a $\Phi (G)\triangleq \bigcap \mathcal{M}_{G}$.

De otro lado, un operador binario común entre elementos de un grupo es el \emph{conmutador}, 
$\left[ a,b \right] |\triangleq a^{-1} \cdot b^{-1} \cdot a \cdot b$. 
Así $\left[{H,K}\right]\triangleq \left\lbrace\left[ h,k \right] 
\vert h  \in H \wedge k \in K\right\rbrace $. 
El conmutador así definido de dos subgrupos no tiene porqué 
ser subgrupo, así se suele considerar $\left\langle{\left[{H,K}\right]}\right\rangle$ 
que es el subgrupo generado por el conmutador. A partir de aquí tenemos dos operadores 
sobre grupos. El primero es el subgrupo \emph{derivado} de diversos órdenes, 
$G^{(0)}\triangleq G$, $\forall i \in \mathds{N}; G^{(i+1)}=
\left\langle\left[ G^{(i)},G^{(i)}\right]\right\rangle$. 
Otro operador importante es $\mathsf{L}^{0}(G)\triangleq G$ y $\forall i \in \mathds{N}\;
\mathsf{L}^{i+1}(G)\triangleq \left\langle\left[{\mathsf{L}^{i},G }\right]\right\rangle$. 
Se dice que un grupo es \emph{nilpotente} si hay una natural $n \in \mathds{N}$ tal que 
$\mathsf{L}^{n}(G)=\mathsf{C}_{1}$. Un grupo es \emph{nilpotente} de clase $n$ tal que 
$\mathsf{L}_{n-1}(G) \neq \mathsf{C}_1$ y $\mathsf{L}_{n}(G)=\mathsf{C}_1$.

Un grupo se dice \emph{abelianizado} del grupo $G$ cuando nos referimos al cociente 
de $G$ por el primer subgrupo derivado $G^{(1)}$, esto es 
$G^{\textsf{ab}}\triangleq \nicefrac{G}{G^{(1)}}$.

Así un grupo es \emph{perfecto} si su abelianizado es trivial, 
$G^{\textsf{ab}} = \mathsf{C}_1$, o dicho de otra manera, 
$G = G^{(1)}$.

Un grupo $G$ se dice \emph{cuasisimple} si es perfecto, $G = G^{(1)}$, y 
$\nicefrac{G}{\mathsf{Z}(G)}$ es simple. Hay que entender que 
$\nicefrac{G}{G^{(1)}} \cong Inn(G) \subseteq Aut(G)$, luego la condición 
anterior equivale a que $Inn(G)$ sea simple.

El subgrupo de Fitting de $G$, $\mathsf{F}(G)$, se define como 
$\mathsf{F}(G) \triangleq \left\langle{\left\lbrace{H \in \mathcal{N}_G \vert 
H \textsf{ nilpotente}}\right\rbrace}\right\rangle$.

Ahora definiremos un grupo \emph{semisimple}. Para esto tiene que existir 
$r \in \mathds{N}\setminus\lbrace 0 \rbrace$ tal que existe 
$S\triangleq \lbrace{G_i \in \mathcal{S}_G \vert 
\forall i\; 0<i<r \Rightarrow G_i \textsf{ cuasisimple}}\rbrace$, 
y $\langle S \rangle = G$. Dicho en palabras, un grupo es semisimple si existe 
una subfamilia finita de subgrupos cuasisimples que lo genere.

Un grupo es \emph{soluble} si existe una subfamilia de $\mathcal{N}_G$ ordenada 
como cadena (por ejemplo creciente) por la inclusión, tal que ya no quepan 
subgrupos intermedios (normales de $G$). Esta cadena iría desde 
$\mathsf{C}_{1}$ hasta $G$. Digamos que es 
$\mathsf{C}_{1}\triangleq H_0 \preccurlyeq H_1 \preccurlyeq \ldots 
\preccurlyeq H_{n-1} \preccurlyeq H_n \triangleq G$, de tal forma que 
$\forall i \in \mathds{N}\;0 \neq i\leq n \Rightarrow 
\nicefrac{H_i}{H_{i-1}} \textsf{ es abeliano}$.

Un subgrupo $H$ de $G$ es \emph{subnormal} si existe una subfamilia de $\mathcal{S}_{G}$ 
finita y ordenada en una cadena por inclusión, de tipo ascendente, tal que 
$H\triangleq H_0 \lhd H_1 \lhd \ldots \lhd H_{n-1} \lhd H_n \triangleq G$. 
Es subnormal de clase $n$.

Un subgrupo $H$ de $G$ se dice \emph{componente} si es subnormal y cuasisimple.

Llamamos ``layer'' de un grupo $G$ al producto (que conmuta) de sus componentes. 
También podemos definirlo como el subgrupo normal semisimple más grande que hay en $G$. 
Digo más grande y no maximal porque es único. Se suele denominar $\mathsf{E}(G)$.

El subgrupo de Fitting generalizado es 
$\mathsf{F}^{*}(G) \triangleq \mathsf{F}(G) \mathsf{E}(G)$.

Sobre el cardinal del grupo finito $G$, lo denominamos por 
$\textsf{card }G = \left| G \right| = \mathsf{o}(G)$. 
Para los subgrupos utilizamos la misma denominación, y hablaremos 
habitualmente de orden de un grupo. Para un grupo $H$ cualquiera 
de $G$, hablaremos del índice de un grupo cuando queramos decir 
el cardinal de conjunto de las clases de equivalencia lateral inducidas, 
$\left[ G : H \right] = \left| \nicefrac{G}{\mathcal{R}_H} \right| =
\left| \nicefrac{G}{\mathcal{R}^H} \right| $.

Ahora tenemos el teorema de Lagrange, que nos dice que si $G$ es finito, 
$H   \preccurlyeq    G$ entonces ${\mathsf{o}(H)} \,\vert\, {{\mathsf{o}}(G)}$. 
Esto nos da en principio un punto de conexión con el teorema fundamental de la aritmética.

De esta forma se hacen importantes los \emph{p-grupos} y los 
\emph{p-subgrupos}, que son en general aquellos grupos finitos cuyo 
orden o cardinal es $p^n$ donde $p$ es un número primo y 
$n \in \mathds{N}\setminus\lbrace 0\rbrace$.

Llamamos ${\mathcal{S}\text{yl}}_{p}(G) \triangleq \left\lbrace H \in \mathcal{S}_{G} 
\vert \left| H \right|=p^n \, \wedge \, p^{n}∣\left| G \right| \, \wedge \, p^{n+1}
\nmid \left| G \right| \right\rbrace$.

Así llamamos \emph{p-núcleo} de $G$ a 
$\mathsf{O}_{p}(G) \triangleq \bigcap {\mathcal{S}\text{yl}}_{p}(G)$.

Emplearemos además los teoremas de Syllow, el teorema de Cauchy, el de Lagrange y 
demás hechos que iremos nombrando, aparte de los habituales sobre productos directos, 
semidirectos y otros que sirven para extender grupos, ya nombrados.

\chapter{Los grupos trivial y cíclicos simples}

Partiremos de los que hemos llamado grupos cíclicos (que incluyen el trivial), 
pero considerando solo los de orden primo. Sea 
$\mathds{P} \triangleq {\left\lbrace{{p    \in \mathds{N}}∣{p>1\,   
\wedge   
\,\forall n   
\in \mathds{N}\setminus{\left\lbrace{0,1,p}\right\rbrace}\; n\nmid p}}\right\rbrace}$.

Para comenzar, los grupos de un orden primo $p$ son todos isomorfos entre sí. 
Sean G y H grupos de orden $p   \in \mathds{P}$. La prueba es clara. 
Si $p    \in \mathds{P}$, cualquier subgrupo 
$G^\prime    \preccurlyeq    G$ y $H^\prime    \preccurlyeq    H$ 
tendrán, por el teorema de Lagrange, orden $1$ o $p$. El subgrupo de orden $1$ 
solo puede ser 
$G^\prime = \left\lbrace 1\right\rbrace    \cong    \mathsf{C}_1$, 
o bien $G^\prime = \left\lbrace g   \in G∣g   \neq    1
\right\rbrace    \cong    \mathsf{C}_p$. Idénticamente en 
$H$, $H^\prime = \left\lbrace 1\right\rbrace    \cong    \mathsf{C}_1$, 
o bien $H^\prime = \left\lbrace h   \in G∣h   \neq    1\right\rbrace    
\cong    \mathsf{C}_p$. Como el orden de $G$ y de $H$ es $p>1$, 
$   \exists    g   \in G\setminus\{1\}   \wedge       
\exists    g   \in G\setminus\{1\}$. Forzosamente esos $g$ y $h$ tienen orden $p$, 
ya que si tuviesen orden $1$ serían $g=1$ y $h=1$, que por elección no es así. 
La aplicación, siempre que $g   \in G\setminus\{1\}$ y que $h \in H\setminus\{1\}$, 
$\varphi \,:\,G \longrightarrow H\,::\, g^n \mapsto h^n \; n \in \mathds{N}$ es 
isomorfismo de grupos.

De esta forma los $\mathsf{C}_{p}\;p \in \mathds{P}$, son los grupos que de forma 
inicial es claro que no tienen estructura interna alguna, esto es, 
$\mathcal{S}_{G}{(\mathsf{C}_{p})} = \mathcal{N}_{G}{(\mathsf{C}_{p})} = 
\mathcal{C}_{G}{(\mathsf{C}_{p})} = \left\lbrace{\{1\},G}\right\rbrace$. 
No tienen composición alguna. Aunque posteriormente encontremos grupos simples 
no cíclicos, serán no abelianos y tienen cierta estructura interna ya que por el 
teorema de Cauchy, existirá al menos un elemento de orden $p \in \mathds{P}$ 
para cada $p∣\left|G\right|$ y por lo tanto tendrá subgrupos de esos órdenes 
(de hecho subgrupos cíclicos de orden primo, para cada primo que divida al orden).

En la clasificación que nos ocupa, esto nos da que salvo isomorfismo de grupos, 
todos los grupos de orden $p   \in \mathds{P}$ y también los de orden $1$, son 
únicos. Además como todos los grupos cíclicos son abelianos (esta característica 
es propia de la suma de exponentes naturales con los que se pueden representar 
estos grupos, suma módulo el orden del grupo cíclico).

Para cada grupo cíclico, el grupo de automorfismos internos es trivial, 
y en general $Aut(\mathsf{C}_p)   \cong    \mathsf{C}_{p-1}$. Es fácil ver que 
para orden $p^n , p   \in \mathds{P} , n   \in \mathds{N}\setminus\{0\}$, 
$Aut(\mathsf{C}_{p^n})   \cong    \mathsf{C}_{p^{n-1}    \cdot (p-1)}$. 
Esta función que asigna a cada primo su predecesor, a cada potencia de un primo 
le asigna el producto de la potencia anterior por el predecesor del primo, y que 
además preserva la multiplicación cuando los factores son coprimos, es la función 
$\phi_{E}$ de Euler, que nos da la cantidad de divisores de un número dado, 
menores que ese número (que es el argumento). Así, si 
$n \in \mathds{N}\setminus\{0\}$, por el teorema fundamental de la aritmética, 
existen una cantidad finita de primos, 
$\left\lbrace p_1 , p_2 , \ldots , p_r \right\rbrace    \subset \mathds{P}$ con 
$r   \in \mathds{N}\,   \wedge   \, r>1$ y una cantidad similar de naturales 
estrictamente positivos, $\left\lbrace n_1, n_2 , \ldots , n_r \right\rbrace 
\subset \mathds{N}\setminus\{0\}$ tales que $n = {\prod}_{k=1}^{r} p_k^{n_k}$. 
Así $Aut(\mathsf{C}_n)   \cong    \mathsf{C}_{{\phi_E}(n)}$.

Para cualquier grupo cíclico de orden primo $p$, el diagrama de Hasse de 
los subgrupos es:

\begin{center}
\begin{tikzpicture}
    \label{tab:GruposCíclicosSimples}
    % First, locate each of the nodes and name them
    \node (top) at (0,0) {$\mathsf{C}_p$};
    \node (bottom) [below of=top] {$\{1_{\mathsf{C}_p}\}$};
    % Now draw the lines:
    \draw [black, thick] (top) -- (bottom);
\end{tikzpicture}
\end{center}

Además tenemos una tabla de grupos y órdenes con algunas propiedades:

\begin{table}[h!]
  \begin{center}
    \caption{Grupos cíclicos de orden primo o trivial}
    \label{tab:GruposCíclicosSimples}
    \begin{tabular}{|c|c|c|c|c|c|} % <-- Alignments: 1st column left, 2nd middle and 3rd right, with vertical lines in between
    \hline
      \textbf{} & \textbf{número} &\textbf{numeración} & \textbf{Nombre} & \textbf{} & \textbf{}\\
      \textbf{Orden} & \textbf{de} &\textbf{dentro} & \textbf{del} & \textbf{absolu-} & \textbf{trivial}\\
      \textbf{Orden} & \textbf{grupos} &\textbf{del} & \textbf{del} & \textbf{tamente} & \textbf{}\\
      \textbf{} & \textbf{no} &\textbf{orden} & \textbf{grupo} & \textbf{simple} & \textbf{}\\
      \textbf{} & \textbf{isomorfos} &\textbf{} & \textbf{} & \textbf{} & \textbf{}\\
      \hline\\
      \hline
      1 & 1 & 1 & $\textsf{C}_1$ & Si & Si\\
      \hline
      2 & 1 & 1  & $\textsf{C}_2$ & Si & No\\
      \hline
      3 & 1 & 1  & $\textsf{C}_3$ & Si & No\\
      \hline
      5 & 1 & 1  & $\textsf{C}_5$ & Si & No\\
      \hline
      7 & 1 & 1  & $\textsf{C}_7$ & Si & No\\
      \hline
      11 & 1 & 1  & $\textsf{C}_{11}$ & Si & No\\
      \hline
      13 & 1 & 1  & $\textsf{C}_{13}$ & Si & No\\
      \hline
      17 & 1 & 1  & $\textsf{C}_{17}$ & Si & No\\
      \hline
      19 & 1 & 1  & $\textsf{C}_{19}$ & Si & No\\
      \hline
      23 & 1 & 1  & $\textsf{C}_{23}$ & Si & No\\
      \hline
      29 & 1 & 1  & $\textsf{C}_{29}$ & Si & No\\
      \hline
      31 & 1 & 1  & $\textsf{C}_{31}$ & Si & No\\
      \hline
      37 & 1 & 1  & $\textsf{C}_{37}$ & Si & No\\
      \hline
      39 & 1 & 1  & $\textsf{C}_{39}$ & Si & No\\
      \hline
      41 & 1 & 1  & $\textsf{C}_{41}$ & Si & No\\
      \hline
      43 & 1 & 1  & $\textsf{C}_{43}$ & Si & No\\
      \hline
      47 & 1 & 1  & $\textsf{C}_{47}$ & Si & No\\
      \hline
      53 & 1 & 1  & $\textsf{C}_{53}$ & Si & No\\
      \hline
      59 & 1 & 1  & $\textsf{C}_{59}$ & Si & No\\
      \hline
      61 & 1 & 1  & $\textsf{C}_{61}$ & Si & No\\
      \hline
    \end{tabular}
  \end{center}
\end{table}

\chapter{Primeros grupos compuestos: orden 4}

Ahora nos ocupa el orden $4$. Directamente caben al menos dos posibilidades, $\textsf{C}_{2^2}$ que tiene un solo grupo normal de orden $2$, un elemento de orden $4$ y que es evidentemente abeliano por ser cíclico, y el grupo $\textsf{C}_{2}  \times  \textsf{C}_{2}$ que es también abeliano por serlo cada subgrupo que forma el producto. Además, tiene $3$ grupos normales de orden $2$ y ningún elemento de orden $4$.

En cuanto a si puede haber más grupos no isomorfos a los anteriores, tenemos que ver que, o bien tiene elemento de orden $4$ o no, y necesariamente, por Cauchy, tienen elementos de orden $2$. Si tiene elemento de orden $4$ solo puede ser el grupo cíclico, por definición de grupo cíclico, y si no, ha de tener más de un elemento de orden $2$. Y tienen que ser exactamente $3$. Si tuviese más elementos (necesariamente de orden $2$), tendríamos más de $4$ elementos y el orden ya no sería $4$. En caso de ensayar alguna forma no abeliana con elementos de orden $2$, tomando $2$ generadores (menos sería cíclico), obtenemos siempre más elementos que en el caso abeliano, por lo que incluso con solo $2$ generadores solo cabe el caso abeliano.

El diagrama de Hasse para el grupo cíclico:

\begin{center}
\begin{tikzpicture}
    % First, locate each of the nodes and name them
    \node (arriba) at (0,0) {$\mathsf{C}_{2^2}=\{1,x,x^2,x^3\}$};
    \node[below of=arriba] (medio) {$\mathsf{C}_{2}  \cong  \{1,x^2\}$};
    \node[below of=medio] (abajo) {$\mathsf{C}_{1}  \cong  \{1_{\textsf{C}_{2^2}}\}$};

    % Now draw the lines:
    \draw [black, thick] (arriba) -- (medio);
    \draw [black, thick] (medio) -- (abajo);
\end{tikzpicture}
\end{center}

Y para el producto:

\begin{center}
\begin{tikzpicture}
    % First, locate each of the nodes and name them
    \node (arriba) at (0,0) {$\mathsf{C}_{2}   \times   \mathsf{C}_{2}   \cong   \{1,a,b,a  \cdot b=b  \cdot a\}$};
    \node (mediocentro) at (0,-2) {$\mathsf{C}_{2}  \cong  \{1,a  \cdot b\}$};
    \node (medioderecha) at (-3,-2) {$\mathsf{C}_{2}  \cong  \{1,a\}$};
    \node (medioizquierda) at (3,-2) {$\mathsf{C}_{2}  \cong  \{1,b\}$};
    \node (abajo) at (0,-4) {$\mathsf{C}_{1}  \cong  \{1_{\mathsf{C}_{2}   \times   \mathsf{C}_{2}}\}$};
    % Now draw the lines:
    \draw [black, thick] (arriba) -- (medioderecha);
    \draw [black, thick] (arriba) -- (mediocentro);
    \draw [black, thick] (arriba) -- (medioizquierda);
    \draw [black, thick] (medioderecha) -- (abajo);
    \draw [black, thick] (mediocentro) -- (abajo);
    \draw [black, thick] (medioizquierda) -- (abajo);
\end{tikzpicture}
\end{center}


\begin{table}[h!]
  \begin{center}
    \caption{Grupos de orden $4$}
    \label{tab:GruposOrden4}
    \begin{tabular}{c|c|c|c|c|c} % <-- Alignments: 1st column left, 2nd middle and 3rd right, with vertical lines in between
      \textbf{} & \textbf{número} &\textbf{numeración} & \textbf{Nombre} & \textbf{} & \textbf{}\\
      \textbf{} & \textbf{de} &\textbf{dentro} & \textbf{del} & \textbf{} & \textbf{}\\
      \textbf{orden} & \textbf{grupos} &\textbf{del} & \textbf{grupo} & \textbf{Simple} & \textbf{abeliano}\\
      \textbf{} & \textbf{no} &\textbf{orden} & \textbf{} & \textbf{} & \textbf{}\\
      \textbf{} & \textbf{isomorfos} &\textbf{} & \textbf{} & \textbf{} & \textbf{}\\
      \hline\\
      \hline
      4 & 2 & 1 & $\textsf{C}_4$ & No & Si\\
      \hline
      4 & 2 & 2  & ${\textsf{C}_2}  \times  {\textsf{C}_2}$ & No & Si\\
      \hline
    \end{tabular}
  \end{center}
\end{table}

\chapter{Grupos de orden $2p$ dónde $p \in \mathds{P}_{>2}$}

Colocamos esta clase de grupos en este lugar porque el siguiente orden en el ascenso que nos ocupa es el $6$ y es del tipo dicho, $6=2  \cdot 3$ y $3  \in \mathds{P}_{>2}$ y toda esta clase de órdenes tiene una estructura de grupos muy similar.

Por una parte está el grupo cíclico que siempre existe, que es abeliano, y que además es único por el teorema de clasificación de grupos abelianos finitos. Este grupo tiene $2$ subgrupos normales, característicos, uno de orden $2$ y otro de orden $p$. Así estos órdenes tienen un solo grupo abeliano y es, entonces, necesariamente, el cíclico. Por el teorema de Cauchy tienen un elemento $a$ de orden $2$ y otro $b$ de orden $p$. Entonces $ab$ tendrá orden $mcm(2,p)=2p$ y así es cíclico.

De otro lado están los grupos no abelianos si existiesen. Pero podemos comprobar que el homomorfismo inyectivo $\varphi : \mathsf{C}_2   \longrightarrow   \mathtt{Aut}(\mathsf{C}_p)$ tal que $\varphi (1)=Id_{\mathsf{C}_{p}}$ y siendo $x  \in {\mathsf{C}_{p}}\setminus\{1\}$ el valor $\varphi (x)$ es el automorfismo que asigna $x   \mapsto x^{p-1}$ y de ahí asigna a $x^k \mapsto x^{k(p-1)}$ obtenemos el producto semidirecto $\textsf{C}_{p} {\rtimes}_{\varphi } \mathsf{C}_{2}$. Aquí los $p$ subgrupos de orden $2$ son no-normales, y el de orden $p$ es único, y por lo tanto no solo normal sino característico. Estos grupos se suelen denominar $\mathsf{Dih}_{p} \equiv \mathsf{C}_{p} {\rtimes}_{\varphi } \mathsf{C}_{2}$, y ya tenemos lo necesario para confeccionar su tabla de Cayley. Optaremos por denominar a la imagen del homomorfismo $\varphi_x  \equiv  \varphi (x)$ por razones de mejor legibilidad y escritura más corta.

Sea $(1,a) \in \left(\textsf{C}_{p} \times \textsf{C}_{2}\right)\setminus\left\lbrace{(1,1)}\right\rbrace$, entonces ha de ser $\varphi_{a}(f^\imath)=f^{(p-1)\imath}$, así, $(1,a) \cdot (b^\imath,a^\jmath)=(1 \cdot \varphi_a (b^\imath),a  \cdot a^\jmath)=(b^{(p-1)\imath},a^{\jmath + 1})$. De otro lado, $(b^\imath,a^\jmath) \cdot (1,a)=(b^\imath  \cdot \varphi_a (1),a^\jmath  \cdot a)=(b^{\imath},a^{\jmath + 1})$, y como $(1,a) \cdot (b^\imath,a^\jmath)=(b^{(p-1)\imath},a^{\jmath + 1})  \neq  (b^{\imath},a^{\jmath + 1}) = (b^\imath,a^\jmath) \cdot (1,a)$, al menos para algún valor de $\imath$, $\mathsf{Dih}_{p}$ es necesariamente no-conmutativo.

\[
    (b^i,a^j)*(b^k,a^l) =
    \begin{cases}
        (b^{i+k},a^{l}) & \Leftarrow j = 0 \\
        (b^{i-k},a^{1+l})& \Leftarrow j = 1
    \end{cases}
\]

Y partiendo del resultado anterior hallamos el inverso de $(b^i,a^j)$:

\[
(b^i,a^j)^{-1} =
\begin{cases}
(b^{-i},1) & \Leftarrow j = 0 \\
(b^{i},a)& \Leftarrow j = 1
\end{cases}
\]

Renombrando $\mathfrak{f}\triangleq (b,1)$ y $\mathfrak{g}\triangleq (1,a)$, tenemos una forma de escribir los elementos del grupo de forma más sencilla y corta. Podemos ver que el elemento $(b^i,a^j)=(b^i,1)*(1,a^j)=(b,1)^i * (1,a)^j$ se puede escribir (es igual a) como $\mathfrak{f}^i \mathfrak{g}^j$. De la misma manera $(1,a^j)*(b^i,1)=(\varphi_{a^j}(b^i),a^j)$ y dependiendo del valor de $j$ podremos ponerla como:

\[
(1,a^j)*(b^i,1) = 
\begin{cases}
(b^{i},1) & \Leftarrow j = 0 \\
(b^{-i},a)& \Leftarrow j = 1
\end{cases}
\]
\[
(b^i,a^j)=
(b^i,1)*(1,a^j)=
 \mathfrak{f}^i \mathfrak{g}^j
\]
Y ahora podemos traducir
\[
\mathfrak{g}^j \mathfrak{f}^i = 
\begin{cases}
\mathfrak{f}^i & \Leftarrow j = 0 \\
\mathfrak{f}^{-i} \mathfrak{g}& \Leftarrow j = 1
\end{cases}
\]
En definitiva tenemos una regla de la multiplicación dihédrica:
\[
\mathfrak{g} \mathfrak{f}^i = 
\mathfrak{f}^{-i} \mathfrak{g}
\]
Para los inversos tenemos:
\[
(\mathfrak{f}^i\mathfrak{g}^j)^{-1} =
\begin{cases}
\mathfrak{f}^{-i} & \Leftarrow j = 0 \\
\mathfrak{f}^{i} \mathfrak{g}& \Leftarrow j = 1
\end{cases}
\]

Vemos por lo tanto que $\mathfrak{f}^{i} \mathfrak{g} = \mathfrak{g} \mathfrak{f}^{-i}$ es elemento de orden $2$ para todo $i$.

La correspondiente tabla de Cayley es:
\begin{align*}
\begin{array}{|c|c|c|c|c|c|c|c|}
\hline
1&\mathfrak{f}&\mathfrak{f}^2&\mathfrak{g}&\mathfrak{gf}&\mathfrak{gf}^2\\
\hline
\mathfrak{f}&\mathfrak{f}^2&1&\mathfrak{gf}^2&\mathfrak{g}&\mathfrak{gf}\\
\hline
\mathfrak{f}^2&1&\mathfrak{f}&\mathfrak{gf}&\mathfrak{gf}^2&\mathfrak{g}\\
\hline
\mathfrak{g}&\mathfrak{gf}&\mathfrak{gf}^2&1&\mathfrak{f}&\mathfrak{gf}^2\\
\hline
\mathfrak{gf}&\mathfrak{gf}^2&\mathfrak{g}&\mathfrak{f}^2&1&\mathfrak{f}\\
\hline
\mathfrak{gf}^2&\mathfrak{g}&\mathfrak{gf}&\mathfrak{f}&\mathfrak{f}^2&1\\
\hline
\end{array}
\end{align*}

Podemos ver ahora como es la conjugación de un elemento:

\begin{align}
(\mathfrak{f}^i\mathfrak{g})(\mathfrak{f}^k\mathfrak{g}^l)(\mathfrak{f}^i\mathfrak{g})^{-1} &=&\\
&=\mathfrak{f}^i\mathfrak{g}\mathfrak{f}^k\mathfrak{g}^l(\mathfrak{f}^i\mathfrak{g})^{-1}=&\\
&=\mathfrak{f}^i\mathfrak{f}^{-k}\mathfrak{g}^{1+l}(\mathfrak{f}^i\mathfrak{g})^{-1}=&\\
&= \mathfrak{f}^{i-k}\mathfrak{g}^{1+l}\mathfrak{g} \mathfrak{f}^{-i} =& \\
&= \mathfrak{f}^{i-k}\mathfrak{g}^{2+l} \mathfrak{f}^{-i} =& \\
&= \mathfrak{f}^{i-k}\mathfrak{g}^{l} \mathfrak{f}^{-i} =& \\
&&=\begin{cases}
\mathfrak{f}^{2i-k} \mathfrak{g} & \Leftarrow l = 1\\
\mathfrak{f}^{-k} & \Leftarrow l = 0
\end{cases}
\end{align}

Y desde aquí vemos que $\left\langle \mathfrak f \right\rangle \lhd \mathsf{Dih}_{p}$.

Y como de hecho no hemos utilizado la primalidad de $p$ en ningún sitio, los grupos $\mathsf{Dih}_{n} \triangleq \mathsf{C}_{n}\rtimes_{\varphi }\mathsf{C}_{2}$, $n \in \mathds{N}_{{}>2}$, siendo $\varphi : \mathsf{C}_{2}  \longrightarrow  Aut(\mathsf{C}_{n})$, de forma que $\varphi_1 = Id_{\mathsf{C}_{n}}$ y $\varphi_{a \neq  1}(b^{k \mod n}) = b^{(n-1)k \mod n}$. El orden de $\mathsf{o}(\mathsf{Dih}_{n})=\left|\mathsf{C}_{n} \times  \mathsf{C}_{2}\right|=2n$, y será no-conmutativo.
Podemos ver por cálculo directo en la fórmula general, que los elementos $\mathfrak{g},\mathfrak{gf},\mathfrak{gf}^2, \ldots , \mathfrak{gf}^{n-1}$ son de orden $2$, que son conjugados entre si, y por lo tanto generan subgrupos no-normales y que el subgrupo generado por $\mathfrak{f}$ es característico. La única diferencia entre que sea primo o no, es que cuando no sea primo aparecerán otros subgrupos, que serán subgrupos de $\left\langle\mathfrak{f}\right\rangle$, amén del grupo de automorfismos, que diferirá en su tamaño, ya que si $n$ es primo hay un total de $n-1$ automorfismos de $Aut(\left\langle\mathfrak{f}\right\rangle) = \left\lbrace{\mathfrak{f}\mapsto \mathfrak{f},\mathfrak{f}\mapsto \mathfrak{f}^2,\mathfrak{f}\mapsto \mathfrak{f}^3,\ldots ,\mathfrak{f}\mapsto \mathfrak{f}^{n-1} }\right\rbrace$. De esta forma $\left|Aut(\left\langle\mathfrak{f}\right\rangle)\right|=n-1$ si $n\in\mathds{P}\setminus\{2\}$. De forma más general, el orden será $\phi_{E}(n)$, la función de Euler que cuenta los coprimos con $n$ que son menores que $n$.
Para los elementos de orden $2$, podemos buscar los siguientes isomorfismos: $\left\{{\mathfrak{g}\mapsto\mathfrak{g},\mathfrak{g}\mapsto\mathfrak{g}\mathfrak{f},\mathfrak{g}\mapsto\mathfrak{g}\mathfrak{f}^2,\ldots,,\mathfrak{g}\mapsto\mathfrak{g}\mathfrak{f}^{n-1}}\right\}$ que nos da un total de $n$ elecciones. Así $Aut(\mathsf{Dih}_{n})$ tiene $n\phi_{E}(n)$ elecciones posibles. En el caso que $n=p$ primo, el orden del grupo de automorfismos es $p(p-1)$.

El subgrupo de Frattini de ${\mathsf{Dih}}_{p}$ lo confeccionaremos explicitando primero la familia de maximales:
${\mathcal{M}}_{{\mathsf{Dih}}_{p}} = \left\{ { \langle \mathfrak{f} \rangle,\langle \mathfrak{g} \rangle,\langle
\mathfrak{gf}\rangle,\ldots,\langle \mathfrak{gf}^{p-1}\rangle } \right\} $. Así queda $\Phi_{{\mathsf{Dih}}_{p}} =
\bigcap\mathcal{M}_{{\mathsf{Dih}}_{p}} = \{1\}$.

Además, el centro de este grupo es claramente $\mathsf{Z}\left({{\mathsf{Dih}}_{p}}\right)=\{1\}$, de dónde $Int\left({{\mathsf{Dih}}_{p}}\right)\cong\nicefrac{{{\mathsf{Dih}}_{p}}}{\mathsf{Z}({{\mathsf{Dih}}_{p}})}\cong{{\mathsf{Dih}}_{p}}$, con lo que ordinal de los automorfismos interiores es $2p$, de los $(p-1)^2$ del total de automorfismos.

El subgrupo de Fitting es el generado por los subgrupos normales y nilpotentes, por lo tanto solo tenemos que ver si $\langle \mathfrak{f} \rangle$ es nilpotente. Pero esto es claro porque es un grupo cíclico, y por lo tanto abeliano, por lo que el subgrupo derivado ${\langle \mathfrak{f} \rangle}^{(1)} = \{1\}$. Como no hay otros normales, $\mathsf{F}({\mathsf{Dih}}_{p})=\langle \mathfrak{f} \rangle$.

El subgrupo derivado primero del grupo dihedral es $\left\{ \left[ \mathfrak{g}^{i}\mathfrak{f}^{j},\mathfrak{g}^{k}\mathfrak{f}^{l} \right]=\mathfrak{f}^{2(j-l)} \right\} = \langle \mathfrak{f} \rangle$. Así $\mathsf{Dih}_{p}^{(1)} = \langle \mathfrak{f} \rangle$. Sin embargo ${\mathsf{L}}_{2}\left( {\mathsf{Dih}}_{p}
\right) = {\mathsf{L}}_{1}\left( {\mathsf{Dih}}_{p} \right) = \mathsf{Dih}_{p}^{(1)} = \langle \mathfrak{f} \rangle$.
 Sin embargo ${\mathsf{Dih}}_{p}^{(2)} = \left\{1\right\}$.
 

Para el caso del grupo cíclico $Aut(\mathsf{C}_{2p})\cong \mathsf{C}_{p-1}$.


El diagrama de Hasse de los subgrupos para $\mathsf{C}_{2p}$ será:

\begin{center}
    \begin{tikzpicture}
    % First, locate each of the nodes and name them
    \node (arriba) at (0,0) {$\mathsf{C}_{2p}   \cong   \{1,a,c\triangleq a^2,\ldots,b\triangleq a^p,\ldots,a^{2p-1}\}$};
    \node (medioarriba) at (3,-2) {$\mathsf{C}_{p}  \cong  \{1,c,\ldots,c^{p-1}\}$};
    \node (medioabajo) at (-3,-4) {$\mathsf{C}_{2}  \cong  \{1,b\}$};
    \node (abajo) at (0,-6) {$\mathsf{C}_{1}  \cong  \{1_{\mathsf{C}_{2p}}\}$};
    
    % Now draw the lines:
    \draw [black, thick] (arriba) -- (medioarriba);
    \draw [black, thick] (arriba) -- (medioabajo);
    \draw [black, thick] (medioarriba) -- (abajo);
    \draw [black, thick] (medioabajo) -- (abajo);
    \end{tikzpicture}
\end{center}

El diagrama de Hasse de los subgrupos para $\mathsf{Dih}_{p}$ será:


\begin{center}
    \begin{tikzpicture}
    \label{tab:GruposOrden2p}
    % First, locate each of the nodes and name them
    \node (arriba) at (0,0) {$\mathsf{Dih}_{p}   \cong   \left\langle \mathfrak{g},\mathfrak{f} \right\rangle$};
    \node (medioarriba) at (2,-2) {$\mathsf{C}_{p}  \cong  \{1,\mathfrak{f},\ldots,{\mathfrak{f}}^{p-1}\}$};
    \node (medioabajo1) at (-2,-4) {$\cong\{1,\mathfrak{g}\}$};
    \node (medioabajo2) at (-4,-4) {$\cong\{1,\mathfrak{gf}\}$};
    \node (medioabajo3) at (-6,-4) {$\ldots$};
    \node (medioabajo4) at (-8,-4) {$\mathsf{C}_{2}  \cong  \{1,{\mathfrak{gf}}^{p-1}\}\cong$};
    \node (abajo) at (0,-6) {$\mathsf{C}_{1}  \cong  \{1_{\mathsf{Dih}_{p}}\}$};
    
    % Now draw the lines:
    \draw [black, thick] (arriba) -- (medioarriba);
    \draw [blue, thick] (arriba) -- (medioabajo1);
    \draw [blue, thick] (arriba) -- (medioabajo2);
    \draw [blue, thick] (arriba) -- (medioabajo4);
    \draw [black, thick] (medioarriba) -- (abajo);
    \draw [black, thick] (medioabajo1) -- (abajo);
    \draw [black, thick] (medioabajo2) -- (abajo);    
    \draw [black, thick] (medioabajo4) -- (abajo);
    \end{tikzpicture}
\end{center}





La tabla con la que ya podemos aumentar las tablas de grupos por órdenes es:





\begin{table}[h!]
    \begin{center}
        \caption{Grupos de orden $2p$}
        \label{tab:GruposOrden2p}
        \begin{tabular}{c|c|c|c|c|c|c} % <-- Alignments: 1st column left, 2nd middle and 3rd right, with vertical lines in between
            \textbf{} & \tiny\textbf{número} &\tiny\textbf{numeración} & \tiny\textbf{nombre} & \textbf{} & {} & {} \\
            \textbf{} & \tiny\textbf{de} &\tiny\textbf{dentro} & \tiny\textbf{del} & \tiny\textbf{} & {} & {} \\
            \tiny\textbf{orden} & \tiny\textbf{grupos} & \tiny\textbf{del} & \tiny\textbf{grupo} & \tiny\textbf{abeliano}  & \tiny\textbf{nilpotente} & \tiny\textbf{resoluble} \\
            \textbf{} & \tiny\textbf{no} & \tiny\textbf{orden} & \textbf{} & \textbf{} & {} & {} \\
            \textbf{} & \tiny\textbf{isomorfos} &\textbf{} & \textbf{} & \textbf{} & {} & {} \\
            \hline
            6 & 2 & 1 & ${\textsf{C}}_{2*3}$ & Si & Si & Si \\
            \hline
            6 & 2 & 2  & ${{\textsf{Dih}}_{3}}  \cong  {{\textsf{C}}_{3}}\rtimes{{{\textsf{C}}_{2}}}$ & No & No & Si \\
            \hline
            10 & 2 & 1 & ${\textsf{C}}_{2*5}$ & Si & Si & Si \\
            \hline
            10 & 2 & 2  & ${{\textsf{Dih}}_{5}}  \cong  {{\textsf{C}}_{5}}\rtimes{{{\textsf{C}}_{2}}}$ & No & No & Si \\
            \hline
            14 & 2 & 1 & ${\textsf{C}}_{2*7}$ & Si & Si & Si \\
            \hline
            14 & 2 & 2  & ${{\textsf{Dih}}_{7}}  \cong  {{\textsf{C}}_{7}}\rtimes{{{\textsf{C}}_{2}}}$ & No & No & Si \\
            \hline
            22 & 2 & 1 & ${\textsf{C}}_{2*11}$ & Si & Si & Si \\
            \hline
            22 & 2 & 2  & ${{\textsf{Dih}}_{11}}  \cong  {{\textsf{C}}_{11}}\rtimes{{{\textsf{C}}_{2}}}$ & No & No & Si \\
            \hline
            26 & 2 & 1 & ${\textsf{C}}_{2*13}$ & Si & Si & Si \\
            \hline
            26 & 2 & 2  & ${{\textsf{Dih}}_{13}}  \cong  {{\textsf{C}}_{13}}\rtimes{{{\textsf{C}}_{2}}}$ & No & No & Si \\
            \hline
            34 & 2 & 1 & ${\textsf{C}}_{2*17}$ & Si & Si & Si \\
            \hline
            34 & 2 & 2  & ${{\textsf{Dih}}_{17}}  \cong  {{\textsf{C}}_{17}}\rtimes{{{\textsf{C}}_{2}}}$ & No & No & Si \\
            \hline
            38 & 2 & 1 & ${\textsf{C}}_{2*19}$ & Si & Si & Si \\
            \hline
            38 & 2 & 2  & ${{\textsf{Dih}}_{19}}  \cong  {{\textsf{C}}_{19}}\rtimes{{{\textsf{C}}_{2}}}$ & No & No & Si \\
            \hline
            46 & 2 & 1 & ${\textsf{C}}_{2*23}$ & Si & Si & Si \\
            \hline
            46 & 2 & 2  & ${{\textsf{Dih}}_{23}}  \cong  {{\textsf{C}}_{23}}\rtimes{{{\textsf{C}}_{2}}}$ & No & No & Si \\
            \hline
            58 & 2 & 1 & ${\textsf{C}}_{2*29}$ & Si & Si & Si \\
            \hline
            58 & 2 & 2  & ${{\textsf{Dih}}_{29}}  \cong  {{\textsf{C}}_{29}}\rtimes{{{\textsf{C}}_{2}}}$ & No & No & Si \\
            \hline
            62 & 2 & 1 & ${\textsf{C}}_{2*31}$ & Si & Si & Si \\
            \hline
            62 & 2 & 2  & ${{\textsf{Dih}}_{31}}  \cong  {{\textsf{C}}_{31}}\rtimes{{{\textsf{C}}_{2}}}$ & No & No & Si \\
            \hline
        \end{tabular}
    \end{center}
\end{table}








    Queda aún una pregunta por responder: ¿Hay más grupos que los dichos del orden $2p$?
    
    De tipo abeliano claramente no, ya se dio la explicación, pero también se puede aludir al teorema de clasificación de los grupos abelianos y finitos.
    
    Para el caso no-abeliano, no cabe que existan más subgrupos normales, pues en ese caso el grupo pasaría a ser abeliano por ser ambas partes del producto abelianas. 
    
    Pensando en p-subgrupos y p-subgrupos de Syllow, siempre tendremos $2$ familias de subgrupos de Syllow, $\mathcal{S}yl_2 (G)$ y $\mathcal{S}yl_p (G)$. Los p-subgrupos serán necesariamente normales porque $[G:H_p]=2$, de dónde los elementos de $\mathcal{S}yl_2 (G)$ serán no-normales. Así que lo que queda por dilucidar es si puede ser $\left|\mathcal{S}yl_{p}(G)\right|>1$. Pongamos que sea así: entonces existirán $H_p^1,H_p^2$ p-subgrupos. Como son diferentes, por hipótesis, $H_p^1 \cap H_p^2 = \{1\}$, y así, si hay $r$ p-subgrupos, tendremos un total de $r(p-1)+1$ elementos debido a los p-subgrupos, y $r\le 2$, de dónde $r=2$ significa que solo habría un 2-subgrupo, pero entonces el 2-subgrupo habría de ser normal y de ahí el grupo abeliano. Luego no hay más grupos no-abelianos que ${\mathsf{Dih}}_{p}$.
    
    \chapter{Grupos de orden $8$}
    
    Los grupos de orden $8$ podrían ser tratados como grupos de orden $p^3$ con $p\in\mathds{P}$, de la misma forma que los grupos de orden $4$ podrían haber sido tratados como los grupos de orden $p^2$. No lo hago así por algunas diferencias entre lo grupos con orden potencia de un primo impar y los grupos con orden potencia de $2$.
    
    Los grupos abelianos de orden $8$ son fáciles de enumerar gracias al teorema de clasificación de los grupos finitos abelianos. Salvo isomorfismo solo pueden ser $\mathsf{C}_{2}\times\mathsf{C}_{2}\times\mathsf{C}_{2}$, $\mathsf{C}_{4}\times\mathsf{C}_{2}$ y $\mathsf{C}_{8}$. Al ser abelianos y por su construcción mediante producto directo podemos reconstruir fácilmente su tablas de Cayley, su retículo de subgrupos (se invierte el teorema de Lagrange) y todos son normales.
    
    Los diagramas de Hasse correpondientes son:
    
    
    
    
    \begin{flushleft}
        \begin{tikzpicture}
        % First, locate each of the nodes and name them
        \node (C8) at (3,0) {$\mathsf{C}_{8}\cong\langle a\vert a^8\rangle$};
        \node (C4xC2) at (-1,0) {$\mathsf{C}_4\times\mathsf{C}_2 \cong  \left\langle a,b \vert a^4=b^2 \right\rangle$};
        \node (C2xC2xC2) at (-7,0) {$\mathsf{C}_2 \times\mathsf{C}_2 \times\mathsf{C}_2 \cong\left\langle a,b,c \vert a^2=b^2=c^2\right\rangle $};
        \node (C8C4) at (3,-2) {$\cong\left\langle a^2\right\rangle $};
        \node (C8C4C2) at (3,-4) {$\cong\left\langle{a^4}\right\rangle$};
        \node (C8abajo) at (3,-6) {$\cong\left\langle{1_{\mathsf{C}_8}}\right\rangle$};
        \node (C4xC2C4) at (-2,-2) {$\mathsf{C}_4 \cong\left\langle{a}\right\rangle$};
        \node (C4xC2C4C2) at (-2,-4) {$\cong\left\langle{a^2}\right\rangle$};
        \node (C4xC2C2) at (1,-4) {$\cong\left\langle{b}\right\rangle$};
        \node (C4xC2abajo) at (-1,-6) {$\cong\left\langle{1_{\mathsf{C}_4 \times \mathsf{C}_2}}\right\rangle$};
        \node (C2xC2xC21) at (-5,-2) {$\cong\left\langle{a,b}\right\rangle$};
        \node (C2xC2xC22) at (-7,-2) {$\cong\left\langle{a,c}\right\rangle$};
        \node (C2xC2xC23) at (-9,-2) {$\mathsf{C}_{2}\times\mathsf{C}_{2}\cong\left\langle{b,c}\right\rangle$};
        \node (C2xC2xC212) at (-5,-4) {$\cong\left\langle{a}\right\rangle$};
        \node (C2xC2xC222) at (-7,-4) {$\cong\left\langle{b}\right\rangle$};
        \node (C2xC2xC232) at (-9,-4) {$\mathsf{C}_{2}\cong\left\langle{c}\right\rangle$};
        \node (C2xC2xC2abajo) at (-7,-6) {$\mathsf{C}_{1}\cong\left\langle{1_{\mathsf{C}_2 \times\mathsf{C}_2 \times\mathsf{C}_2}}\right\rangle$};
        % Now draw the lines:
        \draw [black, thick] (C8) -- (C8C4);
        \draw [black, thick] (C8C4) -- (C8C4C2);
        \draw [black, thick] (C8C4C2) -- (C8abajo);
        \draw [black, thick] (C4xC2) -- (C4xC2C4);
        \draw [black, thick] (C4xC2) -- (C4xC2C2);
        \draw [black, thick] (C4xC2C4) -- (C4xC2C4C2);
        \draw [black, thick] (C4xC2C2) -- (C4xC2abajo);
        \draw [black, thick] (C4xC2C4C2) -- (C4xC2abajo);
        \draw [black, thick] (C2xC2xC2) -- (C2xC2xC21);%ab
        \draw [black, thick] (C2xC2xC2) -- (C2xC2xC22);%ac
        \draw [black, thick] (C2xC2xC2) -- (C2xC2xC23);%bc
        \draw [black, thick] (C2xC2xC21) -- (C2xC2xC212);%ab->a
        \draw [black, thick] (C2xC2xC21) -- (C2xC2xC222);%ab->b
        \draw [black, thick] (C2xC2xC22) -- (C2xC2xC212);%ac->a
        \draw [black, thick] (C2xC2xC22) -- (C2xC2xC232);%ac->c
        \draw [black, thick] (C2xC2xC23) -- (C2xC2xC222);%bc->b
        \draw [black, thick] (C2xC2xC23) -- (C2xC2xC232);%bc->c
        \draw [black, thick] (C2xC2xC212) -- (C2xC2xC2abajo);
        \draw [black, thick] (C2xC2xC222) -- (C2xC2xC2abajo);
        \draw [black, thick] (C2xC2xC232) -- (C2xC2xC2abajo);    
        \end{tikzpicture}
    \end{flushleft}
    
     
    
    
    Los grupos no-abelianos han de tener $\mathsf{o}\left(\mathsf{Z}(G)\right) \in \{2,4\}$, por ser un 2-grupo. Si el orden del centro fuese $2$, $\nicefrac{G}{\mathsf{Z}_{G}}\cong \mathsf{C}_{4}$ ó $\nicefrac{G}{\mathsf{Z}_{G}}\cong \mathsf{C}_{2}\times\mathsf{C}_{2}$. Pero el caso $\nicefrac{G}{\mathsf{Z}_{G}}\cong \mathsf{C}_{4}$ hemos de rechazarlo, por que si ese cociente es cíclico entonces $G$ es abeliano contra la hipótesis. Si el orden del centro fuese $4$, $\nicefrac{G}{\mathsf{Z}_{G}}\cong \mathsf{C}_{2}$, sería cíclico y por lo tanto abeliano contra la hipótesis.
    
    Luego $G$ no-abeliano, $\mathsf{G}=8$, entonces $\mathsf{Z}\left(G\right)\cong\mathsf{C}_{2}$, y $\nicefrac{G}{\mathsf{Z}_{G}}\cong \mathsf{C}_{2}\times\mathsf{C}_{2}$.
    
    Ahora falta ver si $\mathcal{S}yl_{2}(G)=\left\lbrace H\preccurlyeq G \vert \mathsf{o}(H)=4 \right\rbrace$ contiene grupos normales o no-normales, pero en cualquier caso isomorfos entre sí. Caben dos posibilidades, que $\left| \mathcal{S}yl_2 (G)\right|=1$ y que $\left| {{\mathcal{S}yl}_{2}}(G)\right|>1$.
    
    Exploremos $\left| {{\mathcal{S}yl}_{2}}(G)\right|>1$:
    
    El primer caso es $\left| {{\mathcal{S}yl}_{2}}(G)\right|=2$, 
tomando que la intersección es $\{1\}$, hacen un total de $2(4-1)+1=7$ elementos 
en los 2-subgrupos de Syllow, de dónde necesitaríamos un elemento extra que sería 
de orden $2$. Esto hace que el número de elementos final se duplicara y excediera 
el orden buscado. Pongamos por lo tanto que la intersección de los 2-subgrupos de 
Syllow es un subgrupo de orden $2$, de dónde $2(4-2)+2=6$, y necesitaríamos de nuevo 
algún elemento extra de orden $2$, que haría que el orden de $G$ fuese mayor que $8$. 
Por lo tanto este caso se puede desechar.
    
    El segundo caso es $\left| {{\mathcal{S}yl}_{2}}(G)\right|=3$, pero $3*4=12 > 8$, 
por lo que habremos de tomar que todos los subgrupos de Syllow tienen una intersección 
común de orden $2$. Salen $3(4-2)+2=8$ elementos exactamente. De aquí sabemos 
que ha de existir en cada subgrupo de Syllow $2$ elementos que sólo están en ese subgrupo. 
Llamemos $ {{\mathcal{S}yl}_{2}}\left(G\right)=\left\{ N_1,N_2,N_3 \right\} $, 
llamemos 
$H = N_1 \cap N_2 = N_1 \cap N_3 = N_2 \cap N_3 = 
N_1 \cap N_2 \cap N_3 = \left\{ \overline{1},1 \right\}$. 
Sea $N_1=\left\{ \overline{a},a\right\}\cup H$, $N_2=\left\{\overline{b},b\right\}\cup H$ y 
$N_3=\left\{ \overline{c},c\right\}\cup H$. Ahora la pregunta es ¿cuales pueden ser los
ordenes de $a$? Pueden ser $2$ o $4$. Si fuese $2$, $K_1 = \{1,a\}$, $K_2 = \{1,\overline{a}\}$, 
y así con $b$ y $c$. Pero entonces todos los subgrupos de orden $2$ serían isomorfos, 
y los subgrupos de Syllow serían el 4-grupo de Klein, pero 
también lo serían $\{1,a\}\times\{1,b\}$ y saldrían muchas más de $3$ grupos de Syllow. 
Así que solo nos queda que los 
grupos de Syllow sean isomorfos a $\mathsf{C}_{4}$. 
Por lo tanto el orden de los elementos $\{a,\overline{a},b,\overline{b},c,\overline{c}\}$ 
será $4$. Sabemos que $a^2\neq\overline{a}$ pues si no su orden ya
sería mayor que $4$. Por lo tanto solo queda 
$a^2=\overline{a}^2=b^2=\overline{b}^2=c^2=\overline{c}^2=\overline{1}$ y 
además $\overline{1}^2=1$. Sólo queda ver cómo son el resto de productos. Por lo pronto
$a\overline{1}=\overline{1}a=\overline{a}$, $\overline{a}\overline{1}=\overline{1}\overline{a}=a$, 
e igualmente con $b$ y $c$. A todo esto ha de quedar ya claro que $a^{-1}=\overline{a}$ 
y así con $b$ y $c$, y que $\overline{1}^{-1}=\overline{1}$.
               
    Queda ver solamente como es el producto $ab$ etc. Sabemos que $ab\neq\overline{1}$ 
porque entonces $b=\overline{a}$, y con el mismo tipo de argumento llegamos a que 
$ab \in \{c,\overline{c}\}$. Haremos la prueba con $ab=c$, $bc=a$, $ca=b$, y por la 
no-conmutatividad $ba=\overline{c}$, $cb=\overline{a}$, $ac=\overline{b}$. 
Si hacemos la prueba con $ab=\overline{c}$, $bc=a$, $ca=b$, y por la no-conmutatividad  
$ba=c$, $cb=a$, $ac=b$, y obtenemos un grupo isomorfo al anterior.
    
    Nos quedamos con $ab=c$, $bc=a$, $ca=b$, $ba=\overline{c}$, $cb=\overline{a}$, 
$ac=\overline{b}$, $a\overline{1}=\overline{1}a=\overline{a}$, 
$\overline{a}\overline{1}=\overline{1}\overline{a}=a$, 
$b\overline{1}=\overline{1}b=\overline{b}$, 
$\overline{b}\overline{1}=\overline{1}\overline{b}=b$, 
$c\overline{1}=\overline{1}c=\overline{c}$, 
$\overline{c}\overline{1}=\overline{1}\overline{c}=c$, 
$a^2=\overline{a}^2=b^2=\overline{b}^2=c^2=\overline{c}^2=1$, 
$a\overline{a}=b\overline{b}=c\overline{c}=\overline{1}^2=1$.
    
    Definimos al grupo del párrafo anterior como $\mathsf{Q}_8 = \langle a,b | a^4=b^4=(ab)^4=1, a^2=b^2=(ab)^2, ab=a^{-1}b^{-1} \rangle$. Es fácil ver que $\langle a\rangle \lhd \mathsf{Q}_8$, $\langle b\rangle \lhd \mathsf{Q}_8$, $\langle ab\rangle \lhd \mathsf{Q}_8$ (son todos los subgrupos de orden $4$) y que $\mathsf{Z}(\mathsf{Q}_8)=\{1,\overline{1}\}\lhd\mathsf{Q}_8$ que es único subgrupo de orden $2$. Es un grupo hamiltoniano (dedekind no-abeliano).
    
    El diagrama de Hasse de $\mathsf{Q}_8$ será:
    
    
    \begin{center}
        \begin{tikzpicture}
        \label{tab:DiagramaDeHasseParaElGrupoQuaternion8}
        % First, locate each of the nodes and name them
        \node (arriba) at (0,0) {$\mathsf{Q}_{8}   \cong   \left\langle a,b \right\rangle$};
        \node (medioarribaderecha) at (2,-2) {$\cong  \left\langle ab \right\rangle$};
        \node (medioarriba) at (0,-2) {$\cong\left\langle b\right\rangle $};
        \node (medioarribaizquierda) at (-2,-2) {$\mathsf{C}_{4}  \cong\left\langle a\right\rangle $};
        \node (medioabajo) at (0,-4) {$\mathsf{C}_{2}\cong\left\langle{a^2=b^2=(ab)^2}\right\rangle$};
        \node (abajo) at (0,-6) {$\mathsf{C}_{1}  \cong  \{1_{\mathsf{Q}_{8}}\}$};
        
        % Now draw the lines:
        \draw [black, thick] (arriba) -- (medioarribaderecha);
        \draw [black, thick] (arriba) -- (medioarriba);
        \draw [black, thick] (arriba) -- (medioarribaizquierda);
        \draw [black, thick] (medioarribaderecha) -- (medioabajo);
        \draw [black, thick] (medioarriba) -- (medioabajo);
        \draw [black, thick] (medioarribaizquierda) -- (medioabajo);    
        \draw [black, thick] (medioabajo) -- (abajo);
        \end{tikzpicture}
    \end{center}
    
    Antes de pasar a los subgurpos característicos y con nombre propio o a los grupos de automorfismos aún podemos ver algo interesante. Veamos como se factoriza (split) el grupo $\mathsf{Q}_{8}$ en una secuencia exacta corta desde $\mathsf{C}_{4}$ hasta $\mathsf{C}_{2}$:



\begin{center}
\begin{align*}
& \bold{1} 
\longrightarrow 
\langle a \rangle 
\underset{{\substack{{a\mapsto a}\\{b\mapsto b}}}}{\overset{\alpha }{\longrightarrow}}
\mathsf{Q}_{8} 
\underset{{\substack{{a\mapsto 1}\\{b\mapsto \overline{1}}}}}{\overset{\beta_a }{\longrightarrow}}
\langle \overline{1} \rangle 
\longrightarrow 
\bold{1}\\
& \bold{1} 
\longrightarrow 
\langle b \rangle 
\underset{{\substack{{a\mapsto a}\\{b\mapsto b}}}}{\overset{\alpha }{\longrightarrow}}
\mathsf{Q}_{8} 
\underset{{\substack{{a\mapsto \overline{1}}\\{b\mapsto 1}}}}{\overset{\beta_b }{\longrightarrow}}
\langle \overline{1} \rangle 
\longrightarrow 
\bold{1} 
\end{align*}
\end{center}    


	De aquí podemos pensar en modelar una acción de $\mathsf{C}_{4}$ en $\mathsf{C}_{2}$ que nos de la operación propia de $\mathsf{Q}_{8}\equiv\mathsf{Dic}_{2}$. Expongo primero la dos acciones, (la operación tiene que variar del producto semidirecto, y resulta haber una segunda acción de $\mathsf{C}_{2} \times \mathsf{C}_{2}$ sobre $\mathsf{C}_{4}$ y que depende de las dos entradas de orden $2$ en el par ordenado, consistiendo la acción por la multiplicación a derechas por un determinado elemento de $\mathsf{C}_{4}$ determinado por la pareja de ${\mathsf{C}_{2}}{\times}{\mathsf{C}_{2}}$): 

\begin{align*}
&\phi \;:\;\mathsf{C}_{2} \longrightarrow \mathtt{Aut}⁡\left(\mathsf{C}_{4}\right)\; ::\; y \mapsto \langle x^k\mapsto x^{4-k}\rangle
\end{align*}
	
	La siguiente función elije un elemento de $\mathsf{C}_{4}$ en función de $\mathsf{C}_{2}\times\mathsf{C}_{2}$ (constituye un homomorfismo): 

\begin{align*}
&\zeta \;:\; 
{{{\mathsf{C}}_{2}}{\times}{{\mathsf{C}}_{2}}}
\longrightarrow
{{\mathsf{C}}_{4}}\;::\;\left(y,z\right) \mapsto
\begin{cases}
b^2 \Leftarrow {y=z\neq 1}\\
1 \Leftarrow (y\neq z)\vee(y=z=1)
\end{cases}
\end{align*}

	Ahora la acción que buscamos es la multiplicación por la derecha del segundo componente del par por el valor devuelto por $\zeta$.

\begin{align*}
&\theta \;:\; 
{{{\mathsf{C}}_{2}}{\times}{{\mathsf{C}}_{2}}{\times}{{\mathsf{C}}_{4}}}
\longrightarrow
{{\mathsf{C}}_{4}}\;::\;\left(y,z,w\right) \mapsto w\zeta\left(y,z\right)
\end{align*}


	Desempaquetada la acción queda:


\begin{align*}
&\theta \;:\; 
{{{\mathsf{C}}_{2}}{\times}{{\mathsf{C}}_{2}}{\times}{{\mathsf{C}}_{4}}}
\longrightarrow
{{\mathsf{C}}_{4}}\;::\;\left(y,z,w\right) \mapsto
\begin{cases}
wf^2 \Leftarrow {y=z\neq 1}\\
w \Leftarrow (y\neq z)\vee(y=z=1)
\end{cases}
\end{align*}


	La operación quedaría:
	
\begin{align*}
&(x_1,y_1),(x_2,y_2)\in \mathsf{C}_{2}\times\mathsf{C}_{4}\\
&\{a,1\}\cong\mathsf{C}_{2}\\
&\{f,f^2,f^3,1\}\cong\mathsf{C}_{4}\\
&(x_1,y_1)⋅(x_2,y_2)=(x_1⋅x_2, \phi(x_2)(y_1)⋅y_2 \zeta(x_1,x_2)⋅)\\
&(x_1,y_1) (x_2,y_2)=
\begin{cases}
&(x_1,y_1 y_2\zeta(x_1,1) \Leftarrow x_2 = 1\\
&(x_1 a,y_1^{-1} y_2\zeta(x_1,a)) \Leftarrow x_2 = a
\end{cases}
\end{align*}

	Hasta el momento, la operación sería la del grupo $\mathsf{Dih}_{4}$ si no fuese por el aditamento de la acción 
$\theta : \mathsf{C}_{2} \times \mathsf{C}_{2}\times \mathsf{C}_{4} \longleftarrow \mathsf{C}_{4}$, que actúa sobre 
la segunda componente del operando derecho. Hagamos ya las sustitución:
	
\begin{align*}
&(x_1,y_1),(x_2,y_2)\in \mathsf{C}_{2}\times\mathsf{C}_{4}\\
&\{a,1\}\cong\mathsf{C}_{2}\\
&\{f,f^2,f^3,1\}\cong\mathsf{C}_{4}\\
&(x_1,y_1) (x_2,y_2)=
\begin{cases}
&(x_1,y_1 y_2 \zeta(x_1,1) \Leftarrow x_2 = 1\\
&(x_1 a,y_1^{-1} y_2 \zeta(x_1,a) \Leftarrow x_2 = a
\end{cases}
\\
&(x_1,y_1) (x_2,y_2)=
\begin{cases}
&(1,y_1 y_2) \Leftarrow \left(x_2 = 1 \wedge x_1 = 1\right)\\
&(a,y_1 y_2) \Leftarrow \left(x_2 \neq x_1 \right)\\
&(1,y_1^{-1} y_2 f^2) \Leftarrow \left(x_2 = a \wedge x_1 = a\right)
\end{cases}
\end{align*}

	Resumiendo, la operación binaria de grupo de $\mathsf{Q}_{8}$ quedaría:

\begin{align*}
&(x_1,y_1),(x_2,y_2)\in \mathsf{C}_{2}\times\mathsf{C}_{4}\\
&\{a,1\}\cong\mathsf{C}_{2}\\
&\{f,f^2,f^3,1\}\cong\mathsf{C}_{4}\\
&(x_1,y_1) (x_2,y_2)=
\begin{cases}
&(1,y_1 y_2 ) \Leftarrow \left(x_2 = x_1 = 1\right)\\
&(a,y_1 y_2 ) \Leftarrow \left(x_2\neq x_1\right)\\
&(1,y_1^{-1} y_2 f^2) \Leftarrow \left(x_2 = x_1 = a\right)
\end{cases}
\end{align*}

	Concretando la ecuación con los valores que pueden tomar las variables $y_1 , y_2$ que no son otra cosa que potencias de $f$ :
	
\begin{align*}
&\{a,1\}\cong\mathsf{C}_{2}\\
&\{f,f^2,f^3,1\}\cong\mathsf{C}_{4}\\
&(x_1,f^k),(x_2,f^l)\in \mathsf{C}_{2}\times\mathsf{C}_{4}\\
&(x_1,f^k) (x_2,f^l)=
\begin{cases}
&(1,f^{k+l}) \Leftarrow \left(x_2 = x_1 = 1\right)\\
&(a,f^{k+l}) \Leftarrow \left(x_2 \neq x_1 \right)\\
&(1,f^{2l-k}) \Leftarrow \left(x_2 = x_1 = a\right)
\end{cases}
\end{align*}
     
	Para la implementación de esta operación con aritmética modular, nos quedamos exclusivamente con los exponentes, consiguiendo así el algoritmo a implementar en el programa informático que nos dará la tabla, los subgrupos y otros tratamientos del grupo:
	
\begin{align*}
&\{a,1\}\cong\mathsf{C}_{2}\\
&\{f,f^2,f^3,1\}\cong\mathsf{C}_{4}\\
&(a^i,f^k),(a^j,f^l)\in \mathsf{C}_{2}\times\mathsf{C}_{4}
\end{align*}


\begin{align*}
&(i,k)(j,l)=
\begin{cases}
(0,k+l) &\Leftarrow {\left(j = i = 0\right)}\\
(1,k+l) &\Leftarrow {\left( (j = 0 \wedge i=1)\vee(j=1 \wedge i=0) \right)}\\
(0,2l-k) &\Leftarrow {\left(j = i = 1\right)}
\end{cases}
\end{align*}


	La tabla de la operación de grupo resultante (quitado la fila y la columna de argumentos), o tabla de Cayley es, una vez cambiados todas las apariciones de $f^2$ por ${-1}$, y si no aparece solo el elemento ${-1}$, por ejemplo en $x{-1}$, cambiamos la expresion por ${-x}$. Hacemos otra trasformación, cambiamos $a$ por $\imath$, $f$ por $\jmath$ y $af$ por $\kappa$, y vemos si se parece más al grupo cuaternión en su forma habitual:


\begin{align*}
\begin{array}{|c|c|c|c|c|c|c|c|}
\hline
1&{\jmath}&{-1}&{-{\jmath}}&{\imath}&{\kappa}&{-{\imath}}&{-{\kappa}}\\
\hline
{\jmath}&{-1}&{-{\jmath}}&1&{-{\kappa}}&{\imath}&{\kappa}&{-{\imath}}\\
\hline
{-1}&{-{\jmath}}&1&{\jmath}&{-{\imath}}&{-{\kappa}}&{\imath}&{\kappa}\\
\hline
{-{\jmath}}&1&{\jmath}&{-1}&{\kappa}&{-{\imath}}&{-{\kappa}}&{\imath}\\
\hline
{\imath}&{\kappa}&{-{\imath}}&{-{\kappa}}&{-1}&{-{\jmath}}&1&{\jmath}\\
\hline
{\kappa}&{-{\imath}}&{-{\kappa}}&{\imath}&{\jmath}&{-1}&{-{\jmath}}&1\\
\hline
{-{\imath}}&{-{\kappa}}&{\imath}&{\kappa}&1&{\jmath}&{-1}&{-{\jmath}}\\
\hline
{-{\kappa}}&{\imath}&{\kappa}&{-{\imath}}&{-{\jmath}}&1&{\jmath}&{-1}\\
\hline
\end{array}
\end{align*}

	\begin{quote}
		Es importante darse cuenta que este método extiende $\mathsf{C}_{4}\times\mathsf{C}_{2}$, de forma disitnta a como lo hace el producto directo y el semidirecto. Así, para los grupos de orden $4n$, aplicando este método obtendremos los grupos $\mathsf{Dic}_{n}$, e incluso podemos pensar que en general tendremos una generalización para los grupos de la forma $p^2 n$.
	\end{quote}	
	
	
		
	Seguimos con el estudio habitual de este grupo.	
     
    Como $\mathsf{C}_{4}\cong\langle a\rangle\cong\langle b\rangle\cong\langle ab\rangle$ son nilpotentes y normales, al igual que $\mathsf{C}_{2}\cong\langle\overline{1}\rangle$, el subgrupo de Fitting es $\mathsf{F}(\mathsf{Q}_8)=\mathsf{Q}_8$, y el grupo de Frattini, $\Phi_{\mathsf{Q}_8}=\langle \overline{1}\rangle$ y coincide con el centro $\Phi_{\mathsf{Q}_8}=\mathsf{Z}(\mathsf{Q}_8)$.
    
    En cuanto a $D_{\mathsf{Q}_8}^{(1)}=\left[{\mathsf{Q}_8,\mathsf{Q}_8}\right]=\langle \overline{1}\rangle$. Por último $D_{\mathsf{Q}_8}^{(2)}=\{1\}$, luego es resoluble.
    
    Para ver si es nilpotente, tenemos $\mathsf{L}_{1}(\mathsf{Q}_8)=D_{\mathsf{Q}_8}^{(1)}=\Phi_{\mathsf{Q}_8}=\mathsf{Z}(\mathsf{Q}_8)=\langle{\overline{1}}\rangle$. Ahora $\mathsf{L}_{2}(\mathsf{Q}_8)=\left[\mathsf{Q}_8,\mathsf{L}_1(\mathsf{Q}_8)\right]=\left[\langle{a,b}\rangle,\langle{\overline{1}}\rangle\right]$. Esto nos da $\mathsf{L}_{2}(\mathsf{Q}_8)=\{1\}$, de dónde es nilpotente.
    
    En cuanto a los automorfismos internos $\mathtt{Int}(\mathsf{Q}_8)\cong \nicefrac{\mathsf{Q}_8}{\mathsf{Z}(\mathsf{Q}_8)}\cong\mathsf{C}_4$. Para los automorfismos en general habría que mantener fijo el centro (es característico), y solo quedan $a\mapsto a,a\mapsto\overline{a},a\mapsto b,a\mapsto\overline{b},a\mapsto {ab},a\mapsto\overline{ab}$ e idénticamente con $b$ en el origen, aunque no todas las elecciones son compatibles: da un máximo de $36$ y tiene que ser menor. Tiene que ser múltiplo de $3$ y de $4$, de dónde sale $\mathsf{o}\left(\mathtt{Aut}(\mathsf{Q}_8)\right)=24$. $\mathtt{Int}(\mathtt{Q}_8)\cong\mathsf{C}_4$ y sin prueba vemos que $\mathtt{Aut}(\mathtt{Q}_8)\cong\mathsf{S}_4$.
    
    Ahora nos quedamos con el caso que queda: $\mathcal{S}yl_2 (G)=\{N\}$. Sabemos que $\mathsf{Z}(G)$ es orden $2$, dónde $\nicefrac{G}{\mathsf{Z}(\mathsf{G})}\cong {{\mathsf{C}}_{2}}\times{{\mathsf{C}}_{2}}$. Sabemos que $N\lhd G$ porque $[G:N]=2$, ya que $\mathsf{o}(N)=4$. Ha de existir un grupo de orden $2$ por Cauchy, pero buscamos alguno que no sea normal en $G$. Sea ese grupo $H=\{1,g\}$ y $N=\{1,f,f^2,f^3\}$, nos proponemos hacer $\mathsf{C}_{4}\rtimes_{\varphi}\mathsf{C}_{2}$ dónde $\forall k \;\varphi_g (f^k)=f^{4-k}$ y $\forall k\;\varphi_1 (f^k)=f^k$. Obtenemos el grupo $G\cong\mathsf{Dih}_{4}$. En este caso queda claro que $K \triangleq\mathsf{Z}\left(\mathsf{Dih}_4\right)={1,f^2}$. Como todos los subgrupos son abelianos todos los subgrupos son nilpotentes. Normales solo serán $N$ y $K$. ${\mathsf{D}}_{\mathsf{Dih}_4}^{(1)} = N$ y ${\mathsf{D}}_{\mathsf{Dih}_4}^{(2)} = \{1\}$ luego es resoluble. Además, $\mathsf{L}_{1}(\mathsf{Dih}_4)=N$, y de aquí $\mathsf{L}_{2}(\mathsf{Dih}_4)=\left[\mathsf{Dih}_4,N\right]=\{1\}$ luego es nilpotente. El subgrupo de Fitting es $\mathsf{F}(\mathsf{Dih}_4)=\langle f\rangle$ y el de Frattini $\Phi_{\mathsf{Dih}_4} = \{1\}$.
    
    La tabla de Cayley es:
    
\begin{align*}
\begin{array}{|c|c|c|c|c|c|c|c|}
\hline
1&f&f^2&f^3&g&gf&gf^2&gf^3\\
\hline
f&f^2&f^3&1&gf^3&g&gf&gf^2\\
\hline
f^2&f^3&1&f&gf^2&gf^3&g&gf\\
\hline
f^3&1&f&f^2&gf&gf^2&gf^3&g\\
\hline
g&gf&gf^2&gf^3&1&f&f^2&f^3\\
\hline
gf&gf^2&gf^3&g&f^3&1&f&f^2\\
\hline
gf^2&gf^3&g&gf&f^2&f^3&1&f\\
\hline
gf^3&g&gf&gf^2&f&f^2&f^3&1\\
\hline
\end{array}
\end{align*}
    
    El diagrama de Hasse de este grupo es:
    
    
    
    \begin{center}
        \begin{tikzpicture}
        % First, locate each of the nodes and name them
        \node (arriba) at (0,0) {$\mathsf{Dih}_{4}   \cong   \left\langle g,f \right\rangle$};
        \node (N) at (4,-2) {$\mathsf{C}_4 \cong  \left\langle f \right\rangle$};
        \node (Z) at (4,-4) {$\cong\left\langle f^2\right\rangle $};
        \node (H) at (0,-4) {$\cong\left\langle g\right\rangle $};
        \node (H1) at (-2,-4) {$\cong\left\langle{gf}\right\rangle$};
        \node (H2) at (-4,-4) {$\cong\left\langle{gf^2}\right\rangle$};
        \node (H3) at (-6,-4) {$\mathsf{C}_{2}\cong\left\langle{gf^3}\right\rangle$};
        \node (abajo) at (0,-6) {$\mathsf{C}_{1}  \cong  \{1_{\mathsf{Dih}_{4}}\}$};
        
        % Now draw the lines:
        \draw [black, thick] (arriba) -- (N);
        \draw [blue, thick] (arriba) -- (H);
        \draw [blue, thick] (arriba) -- (H1);
        \draw [blue, thick] (arriba) -- (H2);
        \draw [blue, thick] (arriba) -- (H3);
        \draw [black, thick] (N) -- (Z);
        \draw [black, thick] (H) -- (abajo);
        \draw [black, thick] (H1) -- (abajo);    
        \draw [black, thick] (H2) -- (abajo);
        \draw [black, thick] (H3) -- (abajo);
        \draw [black, thick] (Z) -- (abajo);
        \end{tikzpicture}
    \end{center}
    
    
    
    




\begin{table}[h!]
    \begin{center}
        \caption{Grupos de orden $8$}
        \label{tab:GruposOrden2p}
        \begin{tabular}{c|c|c|c|c|c|c} % <-- Alignments: 1st column left, 2nd middle and 3rd right, with vertical lines in between
            \textbf{} & \tiny\textbf{número} &\tiny\textbf{numeración} & \tiny\textbf{nombre} & \textbf{} & {} & {} \\
            \textbf{} & \tiny\textbf{de} &\tiny\textbf{dentro} & \tiny\textbf{del} & \tiny\textbf{} & {} & {} \\
            \tiny\textbf{orden} & \tiny\textbf{grupos} & \tiny\textbf{del} & \tiny\textbf{grupo} & \tiny\textbf{abeliano}  & \tiny\textbf{nilpotente} & \tiny\textbf{resoluble} \\
            \textbf{} & \tiny\textbf{no} & \tiny\textbf{orden} & \textbf{} & \textbf{} & {} & {} \\
            \textbf{} & \tiny\textbf{isomorfos} &\textbf{} & \textbf{} & \textbf{} & {} & {} \\
            \hline
            8 & 5 & 1 & ${\textsf{C}}_{8}$ & Si & Si & Si \\
            \hline
            8 & 5 & 2  & ${{\textsf{C}}_{4}}\times{{\textsf{C}}_{2}}$ & Si & Si & Si \\
            \hline
            8 & 5 & 3 & ${{\textsf{C}}_{2}}\times{{\textsf{C}}_{2}}\times{{\textsf{C}}_{2}}$ & Si & Si & Si \\
            \hline
            8 & 5 & 4  & ${{\textsf{Q}}_{8}}$ & No & Si & Si \\
            \hline
            8 & 5 & 5 & ${{\textsf{Dih}}_{4}}  \cong  {{\textsf{C}}_{4}}\rtimes{{{\textsf{C}}_{2}}}$ & No & Si & Si \\
            \hline
        \end{tabular}
    \end{center}
\end{table}



\chapter{Grupos de orden $p^2$ dónde $p \in {\mathds{P}}_{>2}$}


Para abordar los grupos de orden $9$, dado que $9=3^2$, será igualmente fácil tratar el caso general.

Comenzaremos por los grupos abelianos: o bien tiene un elemento $a$ de orden $p^2$ y entonces es cíclico y tenemos $\mathsf{C}_{p^2}\triangleq\langle{a\vert a^{p^2}=1}\rangle$, o bien no tiene ese elemento y tendrá elementos de orden $p$, digamos $b$ y $c$, y tenemos (a la manera del grupo 4 de Klein) $\mathsf{C}_{p}\times\mathsf{C}_{p}\triangleq\langle{b,c\vert b^{p}=c^{p}=1}\rangle$. El diagrama de su retículo de subgrupos es:






\begin{center}
    \begin{tikzpicture}
    % First, locate each of the nodes and name them
    \node (Cp2) at (3,0) {$\mathsf{C}_{p^2}\cong\langle a\vert a^{(p^2)}\rangle$};
    \node (Cp2Cp) at (3,-2) {$\cong\left\langle a^p\right\rangle $};
    \node (Cp2abajo) at (3,-4) {$\cong\left\{1_{\mathsf{C}_{p^2}}\right\}$};
    \node (CpxCp) at (-3,0) {$\mathsf{C}_p\times\mathsf{C}_p \cong  \left\langle b,c \vert b^p=c^p \right\rangle$};
    \node (CpxCpCp1) at (-5,-2) {$\mathsf{C}_p \cong\left\langle{b}\right\rangle$};
    \node (CpxCpCp2) at (-3,-2) {$\cong\left\langle{c}\right\rangle$};
    \node (CpxCpCp3) at (-1,-2) {$\cong\left\langle{bc}\right\rangle$};
    \node (CpxCpabajo) at (-3,-4) {$\textsf{C}_{1}\cong\left\{1_{\mathsf{C}_2 \times \mathsf{C}_2}\right\}$};
% Now draw the lines:
    \draw [black, thick] (Cp2) -- (Cp2Cp);
    \draw [black, thick] (Cp2Cp) -- (Cp2abajo);
    \draw [black, thick] (CpxCp) -- (CpxCpCp1);
    \draw [black, thick] (CpxCp) -- (CpxCpCp2);
    \draw [black, thick] (CpxCp) -- (CpxCpCp3);
    \draw [black, thick] (CpxCpCp1) -- (CpxCpabajo);
    \draw [black, thick] (CpxCpCp2) -- (CpxCpabajo);
    \draw [black, thick] (CpxCpCp3) -- (CpxCpabajo);
    \end{tikzpicture}
\end{center}


A partir de aquí tenemos que buscar los grupos no abelianos. Pero sabemos que al se un p-grupo, $\mathsf{Z}\left(G\right)\neq\{1\}$, y al ser no-abeliano solo queda $\mathsf{Z}\left(G\right)\cong\mathsf{C}_{p}$,
de dónde $\nicefrac{G}{\mathsf{Z}\left(G\right)} \cong \mathsf{C}_{p}$ que es cíclico, y por lo tanto $G$ es abeliano, contra la hipótesis. No existen grupos no abelianos de orden $p^2$.




\begin{table}[h!]
    \begin{center}
        \caption{Grupos de orden $p^2$}
        \label{tab:GruposOrden2p}
        \begin{tabular}{c|c|c|c|c|c|c} % <-- Alignments: 1st column left, 2nd middle and 3rd right, with vertical lines in between
            \textbf{} & \tiny\textbf{número} &\tiny\textbf{numeración} & \tiny\textbf{nombre} & \textbf{} & {} & {} \\
            \textbf{} & \tiny\textbf{de} &\tiny\textbf{dentro} & \tiny\textbf{del} & \tiny\textbf{} & {} & {} \\
            \tiny\textbf{orden} & \tiny\textbf{grupos} & \tiny\textbf{del} & \tiny\textbf{grupo} & \tiny\textbf{abeliano}  & \tiny\textbf{nilpotente} & \tiny\textbf{resoluble} \\
            \textbf{} & \tiny\textbf{no} & \tiny\textbf{orden} & \textbf{} & \textbf{} & {} & {} \\
            \textbf{} & \tiny\textbf{isomorfos} &\textbf{} & \textbf{} & \textbf{} & {} & {} \\
            \hline
            9 & 2 & 1 & ${\textsf{C}}_{9}$ & Si & Si & Si \\
            \hline
            9 & 2 & 2  & ${{\textsf{C}}_{3}}\times{{\textsf{C}}_{3}}$ & Si & Si & Si \\
            \hline
            25 & 2 & 1 & ${{\textsf{C}}_{25}}$ & Si & Si & Si \\
            \hline
            25 & 2 & 2 & ${{\textsf{C}}_{5}}\times{{\textsf{C}}_{5}}$ & Si & Si & Si \\
            \hline
            49 & 2 & 1 & ${{\textsf{C}}_{49}}$ & Si & Si & Si \\
            \hline
            49 & 2 & 2 & ${{\textsf{C}}_{7}}\times{{\textsf{C}}_{7}}$ & Si & Si & Si \\
            \hline
        \end{tabular}
    \end{center}
\end{table}


\chapter{Grupos de orden $12$}

Veremos estos grupos de forma específica, porque el orden $12$ tiene cierta peculiaridad (el grupo alternante) que no tiene análogos en los órdenes $4p$ dónde hay dos familias, $p=1 \mod 4$ y $p=3\mod 4$ (esta última sería la correspondiente al orden $12$).

Primero enumeramos los grupos abelianos de ese orden: $\mathsf{C}_{2}\times\mathsf{C}_{6} \triangleq \langle a,b\vert a^2=b^6=1 \rangle$ y $\mathsf{C}_{12} \triangleq \langle c \vert c^{12}=1 \rangle$.

Sus diagramas de Hasse del retículo de sus subgrupos:



\begin{center}
    \begin{tikzpicture}
    % First, locate each of the nodes and name them
    \node (C12) at (5,15) {$\mathsf{C}_{12}\cong\langle a\vert a^{12}=1\rangle$};
    \node (C12C6) at (12,12) {$\mathsf{C}_{6}\cong\left\langle a^2\right\rangle $};
    \node (C12C4) at (0,9) {$\mathsf{C}_{4}\cong\left\langle a^3\right\rangle $};
    \node (C12C3) at (12,6) {$\mathsf{C}_{3}\cong\left\langle a^4\right\rangle $};
    \node (C12C4C2) at (0,3) {$\mathsf{C}_{2}\cong\left\langle a^6\right\rangle $};
    \node (C12abajo) at (5,0) {$\mathsf{C}_{1}\cong\left\{1_{\mathsf{C}_{12}}\right\}$};
    % Now draw the lines:
    \draw [black, thick] (C12) -- (C12C6);
    \draw [black, thick] (C12) -- (C12C4);
    \draw [black, thick] (C12C6) -- (C12C3);
    \draw [black, thick] (C12C6) -- (C12C4C2);
    \draw [black, thick] (C12C4) -- (C12C4C2);
    \draw [black, thick] (C12C4C2) -- (C12abajo);
    \draw [black, thick] (C12C3) -- (C12abajo);
    \end{tikzpicture}
\end{center}



\begin{center}
    \begin{tikzpicture}
    % First, locate each of the nodes and name them
    \node (a-b) at (6,12)
    {$\mathsf{C}_{6}\times\mathsf{C}_{2}\cong\left\langle{a,b\vert a^6=b^2=1}\right\rangle$};
    \node (a2b) at (3,9) 
    {$\cong \left\langle{a^2 b}\right\rangle$};
    \node (a) at (0,9) 
    {$\mathsf{C}_{6}\cong\left\langle{a}\right\rangle$};
    \node (ab) at (12,9) 
    {$\cong\left\langle{ab}\right\rangle$};
    \node (a3-b) at (9,6) 
    {$\mathsf{C}_{2}\times\mathsf{C}_{2}\cong\left\langle{a^3,b}\right\rangle$};
    \node (a2) at (3,3) 
    {$\mathsf{C}_{3}\cong\left\langle{a^2}\right\rangle$};
    \node (b) at (9,0) 
    {$\cong\left\langle{b}\right\rangle$};
    \node (a3) at (0,0) 
    {$\mathsf{C}_{2}\cong\left\langle{a^3}\right\rangle$};
    \node (a3b) at (12,0) 
    {$\cong\left\langle{a^3 b}\right\rangle$};
    \node (1) at (6,-3) 
    {$\mathsf{C}_{1}\cong\left\langle{1_{\mathsf{C}_{6}\times\mathsf{C}_{2}}}\right\rangle$};
   % Now draw the lines:
    \draw [black, thick] (a-b) -- (ab);
    \draw [black, thick] (a-b) -- (a);
    \draw [black, thick] (a-b) -- (a2b);
    \draw [black, thick] (a-b) -- (a3-b);
    \draw [black, thick] (ab) -- (a3b);
    \draw [black, thick] (ab) -- (a2);
    \draw [black, thick] (a) -- (a2);
    \draw [black, thick] (a) -- (a3);
    \draw [black, thick] (a2b) -- (a2);
    \draw [black, thick] (a2b) -- (b);
    \draw [black, thick] (a3-b) -- (a3b);
    \draw [black, thick] (a3-b) -- (a3);
    \draw [black, thick] (a3-b) -- (b);
    \draw [black, thick] (a2) -- (1);
    \draw [black, thick] (a3b) -- (1);
    \draw [black, thick] (a3) -- (1);
    \draw [black, thick] (b) -- (1);
    \end{tikzpicture}
\end{center}

Ahora comenzaremos con los grupos no abelianos.

Para comenzar veamos las posibilidades que tenemos en cuanto al centro: podría ser $\mathsf{Z}\left(G\right)$ el subgrupo
trivial. También podría ser de órdenes $6$, $4$, $3$ y $2$. Como el centro es abeliano solo puede ser 
$\mathsf{Z}\left(G\right)\cong \mathsf{C}_{6} \;\acute{o}\; \cong \mathsf{C}_{4} \;\acute{o}\; 
\cong \mathsf{C}_{2}\times\mathsf{C}_{2} \;\acute{o}\; \cong \mathsf{C}_{3} \;\acute{o}\; \cong \mathsf{C}_{2} 
\;\acute{o}\; \cong \mathsf{C}_{1}$. Podemos descartar como centro $\mathsf{C}_{6}$ pues el cociente sería 
cíclico, y por lo tanto, $G$ sería abeliano. Podemos utilizar el mismo argumento con los grupos de orden $4$.

 Quedaría las siguientes posibilidades: 
\begin{align}
\mathsf{Z}\left(G\right) &\cong \mathsf{C}_{3} &\text{ si }\quad\quad\quad\quad& 
\nicefrac{G}{\mathsf{Z}\left(G\right)}\cong \mathsf{C}_{2}\times \mathsf{C}_{2}\\
&\cong \mathsf{C}_{2}  &\text{ si }\quad\quad\quad\quad& \nicefrac{G}{\mathsf{Z}\left(G\right)}
\cong \mathsf{C}_{3}\times \mathsf{C}_{2}\\
&\cong \mathsf{C}_{1}&&
\end{align}

Ahora podemos pensar en grupos según tengan un grupos de determinados órdenes normales o no.

Podemos pensar que tenga subgrupos normales de orden $6$, llamémos a uno de ellos $N_6$, o bien que algún grupo de 
orden $6$ sea normal. Este último caso no se puede dar, porque en el caso que hubiese uno solo, digamos $H_6$, entonces $[G:H_6]=2$ y automáticamente $H_6$ sería normal. Luego si hay subgrupos de orden $6$ han de ser normales.

Pero siempre tendremos elementos de orden $3$ y $2$ por el teorema de Cauchy. Subgrupos de orden $4$ existirán por el primer teorema de Syllow. Supondremos que este subgrupo maximal es normal, $N_4$, y el de orden $3$, digamos $H_3$ no lo puede ser, por que sino tenemos que el grupo es abeliano. Estos casos serán del tipo $N_4\rtimes H_3$, o bien a la inversa, y tendríamos $N_3\rtimes H_4$ teniendo en cuenta que los grupos de orden $4$ pueden ser bien $\mathsf{C}_{4}$, bien $\mathsf{C}_{2} \times \mathsf{C}_{2}$.

Enumerando las opciones:
\begin{align}
\mathsf{C}_{4}\rtimes \mathsf{C}_{3}\\
\mathsf{C}_{3}\rtimes \mathsf{C}_{4}\\
\left(\mathsf{C}_{2}\times\mathsf{C}_{2}\right)\rtimes \mathsf{C}_{3}\\
\mathsf{C}_{3}\rtimes \left(\mathsf{C}_{2}\times\mathsf{C}_{2}\right)
\end{align}

En el caso que consideremos que existe grupo de orden $6$ (por lo tanto normal), tendremos las siguientes posibilidades:

\begin{align}
\mathsf{Dih}_{6}\cong\mathsf{C}_{6}\rtimes_{\langle{g\mapsto g\,,\,f\mapsto f^{-1}}\rangle} \mathsf{C}_{2}\\
\mathsf{Dih}_{3} \times \mathsf{C}_{2}\\
\mathsf{Dih}_{3} \rtimes \mathsf{C}_{2}
\end{align}

Para hacer los cálculos necesarios debemos tener los grupos de automorfismos. 

Los automorfismos del grupo dihédrico de orden $6$ son:
\begin{align}
&\texttt{Aut}(\mathsf{Dih}_{3})=\nonumber\\ 
&=\quad
    \left\lbrace \right.
        \varphi_0\triangleq\binom{g\mapsto g}{f\mapsto f}, 
        \varphi_1\triangleq\binom{g\mapsto g}{f\mapsto f^2}, 
        \varphi_2\triangleq\binom{g\mapsto gf}{f\mapsto f},\nonumber\\
&\phantom{=}\quad\varphi_3\triangleq\binom{g\mapsto gf}{f\mapsto f^2}, 			
	\varphi_4\triangleq\binom{g\mapsto gf^2}{f
	\mapsto f}, \varphi_5\triangleq\binom{g\mapsto gf^2}{f\mapsto f^2}
    \left. \right\rbrace\nonumber \cong \mathsf{Dih}_{3}
\end{align}

Los automorfismos del grupo $4$ de Klein son:
\begin{align}
&\texttt{Aut}(\mathsf{C}_{2}\times\mathsf{C}_{2})=\nonumber\\ 
&=\quad
	\left\lbrace \right.
	\phi_0\triangleq\binom{a\mapsto a}{b\mapsto b}, 
	\phi_1\triangleq\binom{a\mapsto a}{b\mapsto ab}, 
	\phi_2\triangleq\binom{a\mapsto ab}{b\mapsto b},\nonumber\\
&\phantom{=}\quad
	\phi_3\triangleq\binom{a\mapsto ab}{b\mapsto a}, 			
	\phi_4\triangleq\binom{a\mapsto b}{b\mapsto ab}, 
	\phi_5\triangleq\binom{a\mapsto b}{b\mapsto a}
	\left. \right\rbrace\nonumber \cong \mathsf{Dih}_{3}
\end{align}

Los automorfismos del grupo $\mathsf{C}_{4}$ son:

\begin{align}
&\texttt{Aut}(\mathsf{C}_{4})=
	\left\lbrace
		\psi_0\triangleq\left({a\mapsto a}\right), 
		\psi_1\triangleq\left({a\mapsto a^3}\right)
	\right\rbrace 
	\cong \mathsf{C}_{2}
\end{align}

Los automorfismos del grupo $\mathsf{C}_{3}$ son:

\begin{align}
&\texttt{Aut}(\mathsf{C}_{3})=
\left\lbrace
	\vartheta_0\triangleq\left({a\mapsto a}\right), 
	\vartheta_1\triangleq\left({a\mapsto a^2}\right)
\right\rbrace 
\cong \mathsf{C}_{2}
\end{align}

Una vez enumerados los distintos grupos de automorfismos plantearemos las acciones 
que se utilizarán en los productos semidirectos, mediante homomorfismos de un grupo 
(el no necesariamente normal) al grupo de automorfismos del otro grupo 
(el que será normal).

Para $\mathsf{C}_{4}\rtimes\mathsf{C}_{3}$ tenemos solo el homomorfismo identidad, pues si 
$\langle a\rangle=\mathsf{C}_{3}$, y teniendo en cuenta $\psi_1$ es de orden 
$2$, $a \mapsto \psi_1$ entonces, por ser homomorfismo, 
$a^2 \mapsto \psi_1\circ\psi_1=\psi_0=Id$. Ahora $a^3=1 \mapsto \psi_1$ 
y llegamos a que no puede ser homomorfismo, ya que esto exigiría $a^3=1 \mapsto \psi_0 = Id$. 
Por lo tanto, por este camino podemos desechar 
$\mathsf{C}_{4}\rtimes_{Id}\mathsf{C}_{3} = 
\mathsf{C}_{4}\times\mathsf{C}_{3}\cong \mathsf{C}_{12}$ 
que es abeliano y ya hemos dado cuenta de él. 

Para $\mathsf{C}_{3}\rtimes\mathsf{C}_{4}$ tenemos una única acción:

\begin{align*}
& \Psi : \mathsf{C}_{4} \longrightarrow \texttt{Aut}\left( \mathsf{C}_{3} \right) \\
& \qquad        1				\mapsto		{\vartheta}_0 = Id \\
& \qquad		a				\mapsto		{\vartheta}_1 \\
& \qquad		a^2				\mapsto		{\vartheta}_0 \\
& \qquad		a^3				\mapsto		{\vartheta}_1
\end{align*}

Así, tomando el conjunto base $\mathsf{C}_{4}\times\mathsf{C}_{3}$ establecemos la 
operación de grupo para $\mathsf{C}_{3}\rtimes_{\Psi} \mathsf{C}_{4}$:

\begin{align*}
& a^4=1 \\
& b^3=1 \\
& (a^i,b^j)(a^k,b^m) \triangleq (a^{i+k} ,{\Psi}_{a^k}(b^j) {b^{m}}) \\
&
(a^i,b^j)(a^k,b^m) = 
	\begin{cases}
		\left( a^{i  } , \vartheta_0(b^{j}){b^m} \right)   &\Leftarrow k=0\\
		\left( a^{i+1} , \vartheta_1(b^{j}){b^m} \right)   &\Leftarrow k=1\\
		\left( a^{i+2} , \vartheta_0(b^{j}){b^m} \right)   &\Leftarrow k=2\\
		\left( a^{i+3} , \vartheta_1(b^{j}){b^m} \right)   &\Leftarrow k=3
	\end{cases} 
\end{align*}

Ahora sustituimos por el valor de $\vartheta_t$:

\begin{align*}
& a^4=1 \\
& b^3=1 \\
& (a^i,b^j)(a^k,b^m) \triangleq (a^{i+k} ,{\Psi}_{a^k}(b^j) {b^{j}}) \\
&
(a^i,b^j)(a^k,b^m) = 
	\begin{cases}
		\left( a^{i  } , {b^{j}} {b^m} \right)   &\Leftarrow k=0\\
		\left( a^{i+1} , {b^{2j}}{b^m} \right)   &\Leftarrow k=1\\
		\left( a^{i+2} , {b^{j}} {b^m} \right)   &\Leftarrow k=2\\
		\left( a^{i+3} , {b^{2j}}{b^m} \right)   &\Leftarrow k=3
	\end{cases} 
\end{align*}


Y obtenemos ya la expresión final:


\begin{align*}
& a^4=1 \\
& b^3=1 \\
& (a^i,b^j)(a^k,b^m) \triangleq (a^{i+k} ,{\Psi}_{a^k}(b^j) {b^{j}}) \\
&
(a^i,b^j)(a^k,b^m) = 
	\begin{cases}
		\left( a^{i  } , {b^{m+j}}  \right)   &\Leftarrow k=0\\
		\left( a^{i+1} , {b^{m+2j}} \right)   &\Leftarrow k=1\\
		\left( a^{i+2} , {b^{m+j}}  \right)   &\Leftarrow k=2\\
		\left( a^{i+3} , {b^{m+2j}} \right)   &\Leftarrow k=3
	\end{cases} 
\end{align*}


En la siguiente tabla de Cayley, en vez de poner $(a^k,b^m)$ 
pondré solo $(k,m)$. La multiplicación cogiendo solo los 
exponentes quedaría así:


\begin{align*}
& a^4=1 \\
& b^3=1 \\
& (a^i,b^j)(a^k,b^m) \triangleq (a^{i+k} ,{\Psi}_{a^k}(b^j) {b^{m}}) \cong\\
& \cong  (i,j)(k,m) \triangleq ({i+k}\mod 4 ,\exp{({\Psi}_{a^k}(b^j))}+m \mod 3) \\
&
(i,j)(k,m) = 
	\begin{cases}
		\left( {i  }\mod 4 , {m+j}  \mod 3 \right)   &\Leftarrow k=0\\
		\left( {i+1}\mod 4 , {m+2j} \mod 3 \right)   &\Leftarrow k=1\\
		\left( {i+2}\mod 4 , {m+j}  \mod 3 \right)   &\Leftarrow k=2\\
		\left( {i+3}\mod 4 , {m+2j} \mod 3 \right)   &\Leftarrow k=3
	\end{cases} 
\end{align*}

Ahora llamando $1 \triangleq (0,0)$, $a^k \triangleq (k,0)$ y $b^m \triangleq (0,m)$, y sabiendo que $a^k b^m = (k,0)(0,m)$:
\\

\begin{flushleft}
\begin{tabular}{|c|c|c|c|c|c|c|c|c|c|c|c|}
\hline
 $1$ & $b$ & $b^2$ & $a$ & $ab$ & $ab^2$ & $a^2$ & $a^2b$ & $a^2b^2$ & $a^3$ & $a^3b$ & $a^3b^2$ \\
\hline
 $b$ & $b^2$ & $1$ & $ab^2$ & $a$ & $ab$ & $a^2b$ & $a^2b^2$ & $a^2$ & $a^3b^2$ & $a^3$ & $a^3b$ \\
\hline
 $b^2$ & $1$ & $b$ & $ab$ & $ab^2$ & $a$ & $a^2b^2$ & $a^2$ & $a^2b$ & $a^3b$ & $a^3b^2$ & $a^3$ \\
\hline
 $a$ & $ab$ & $ab^2$ & $a^2$ & $a^2b$ & $a^2b^2$ & $a^3$ & $a^3b$ & $a^3b^2$ & $1$ & $b$ & $b^2$ \\
\hline
 $ab$ & $ab^2$ & $a$ & $a^2b^2$ & $a^2$ & $a^2b$ & $a^3b$ & $a^3b^2$ & $a^3$ & $b^2$ & $1$ & $b$ \\
\hline
 $ab^2$ & $a$ & $ab$ & $a^2b$ & $a^2b^2$ & $a^2$ & $a^3b^2$ & $a^3$ & $a^3b$ & $b$ & $b^2$ & $1$ \\
\hline
 $a^2$ & $a^2b$ & $a^2b^2$ & $a^3$ & $a^3b$ & $a^3b^2$ & $1$ & $b$ & $b^2$ & $a$ & $ab$ & $ab^2$ \\
\hline
 $a^2b$ & $a^2b^2$ & $a^2$ & $a^3b^2$ & $a^3$ & $a^3b$ & $b$ & $b^2$ & $1$ & $ab^2$ & $a$ & $ab$ \\
\hline
 $a^2b^2$ & $a^2$ & $a^2b$ & $a^3b$ & $a^3b^2$ & $a^3$ & $b^2$ & $1$ & $b$ & $ab$ & $ab^2$ & $a$ \\
\hline
 $a^3$ & $a^3b$ & $a^3b^2$ & $1$ & $b$ & $b^2$ & $a$ & $ab$ & $ab^2$ & $a^2$ & $a^2b$ & $a^2b^2$ \\
\hline
 $a^3b$ & $a^3b^2$ & $a^3$ & $b^2$ & $1$ & $b$ & $ab$ & $ab^2$ & $a$ & $a^2b^2$ & $a^2$ & $a^2b$ \\
\hline
 $a^3b^2$ & $a^3$ & $a^3b$ & $b$ & $b^2$ & $1$ & $ab^2$ & $a$ & $ab$ & $a^2b$ & $a^2b^2$ & $a^2$ \\
\hline
\end{tabular} 
\end{flushleft}

Los órdenes de cada uno de los elementos es:

\begin{align*}
&\mathsf{o}\left(1\right)=1\\
&\mathsf{o}\left(a^2\right)=2\\
&\mathsf{o}\left(b\right)=3\\
&\mathsf{o}\left(b^2\right)=3\\
&\mathsf{o}\left(a\right)=4\\
&\mathsf{o}\left(ab\right)=4\\
&\mathsf{o}\left(ab^2\right)=4\\
&\mathsf{o}\left(a^3\right)=4\\
&\mathsf{o}\left(a^3b\right)=4\\
&\mathsf{o}\left(a^3b^2\right)=4\\
&\mathsf{o}\left(a^2b\right)=6\\
&\mathsf{o}\left(a^2b^2\right)=6
\end{align*}

Por lo pronto podemos ver que se invierte el teorema de Lagrange en este grupo.

Como $a^2$ tiene orden $2$, y es el único elemento de ese orden, es normal luego 
$\langle a^2\rangle\lhd \mathsf{C}_{3}\rtimes_{\Psi} \mathsf{C}_{4}$, 
entonces ha de estar en el centro: 

\[a^2 \in \mathsf{Z}(\mathsf{C}_{3}\rtimes_{\Psi} \mathsf{C}_{4})\]

Además los subgrupos cíclicos son los únicos subgrupos propios (siendo el grupo entero 
el único que no es cíclico) :

\begin{align*}
{\mathcal{L}}\left({{{\mathsf{C}}_{3}}\rtimes{{\mathsf{C}}_{4}}}\right) =
\begin{cases}
&{\left\langle{1}\right\rangle}={\left\lbrace{1}\right\rbrace}\\
&\qquad {\mathsf{o}}{\left({\left\langle{1}\right\rangle}\right)}=1\\
&{\left\langle{a^2}\right\rangle}={\left\lbrace{1,a^2}\right\rbrace}\\
&\qquad {\mathsf{o}}{\left({\left\langle{a^2}\right\rangle}\right)}=2\\
&{\left\langle{b}\right\rangle}={\left\lbrace{1,b,b^2}\right\rbrace}\\
&\qquad {\mathsf{o}}{\left({\left\langle{b}\right\rangle}\right)}=3\\
&{\left\langle{a}\right\rangle}={\left\lbrace{1,a,a^2,a^3}\right\rbrace}\\
&\qquad {\mathsf{o}}{\left({\left\langle{a}\right\rangle}\right)}=4\\
&{\left\langle{a b}\right\rangle}={\left\lbrace{1,ab,a^2,a^3b}\right\rbrace}\\
&\qquad {\mathsf{o}}{\left({\left\langle{a b}\right\rangle}\right)}=4\\
&{\left\langle{a b^2}\right\rangle}={\left\lbrace{1,ab^2,a^2,a^3b^2}\right\rbrace}\\
&\qquad {\mathsf{o}}{\left({\left\langle{a b^2}\right\rangle}\right)}=4\\
&{\left\langle{a^2 b}\right\rangle}={\left\lbrace{1,b,b^2,a^2,a^2b,a^2b^2}\right\rbrace}\\
&\qquad {\mathsf{o}}{\left({\left\langle{a^2 b}\right\rangle}\right)}=6\\
&{\left\langle{a , b}\right\rangle}={{{\mathsf{C}}_{3}}\rtimes{{\mathsf{C}}_{4}}}\\
&\qquad {\mathsf{o}}{\left({\left\langle{a , b}\right\rangle}\right)}=12
\end{cases}
\end{align*}


El subgrupo centro (característico):

\begin{align*}
&\mathsf{Z}\left({{{\mathsf{C}}_{3}}\rtimes{{\mathsf{C}}_{4}}}\right) 
= {\left\langle{a^2}\right\rangle}
\end{align*}

Los subgrupos normalizadores de cada subgrupo en el grupo de estudio, 
nos dará también la información de qué subgrupos son normales:

\begin{align*}
&{\mathsf{N}}_{{{\mathsf{C}}_{3}}\rtimes{{\mathsf{C}}_{4}}}
\left({\left\lbrace{1}\right\rbrace}\right) = 
{{{\mathsf{C}}_{3}}\rtimes{{\mathsf{C}}_{4}}}\\
&{\mathsf{N}}_{{{\mathsf{C}}_{3}}\rtimes{{\mathsf{C}}_{4}}}
\left({\left\langle{a^2}\right\rangle}\right) = 
{{{\mathsf{C}}_{3}}\rtimes{{\mathsf{C}}_{4}}}\\
&{\mathsf{N}}_{{{\mathsf{C}}_{3}}\rtimes{{\mathsf{C}}_{4}}}
\left({\left\lbrace{1,b,b^2}\right\rbrace}\right) = 
{{{\mathsf{C}}_{3}}\rtimes{{\mathsf{C}}_{4}}}\\
&{\mathsf{N}}_{{{\mathsf{C}}_{3}}\rtimes{{\mathsf{C}}_{4}}}
\left({\left\langle{a}\right\rangle}\right) = 
{\left\langle{a}\right\rangle}\\
&{\mathsf{N}}_{{{\mathsf{C}}_{3}}\rtimes{{\mathsf{C}}_{4}}}
\left({\left\langle{ab}\right\rangle}\right) = 
{\left\langle{ab}\right\rangle}\\
&{\mathsf{N}}_{{{\mathsf{C}}_{3}}\rtimes{{\mathsf{C}}_{4}}}
\left({\left\langle{ab^2}\right\rangle}\right) = 
{\left\langle{ab^2}\right\rangle}\\
&{\mathsf{N}}_{{{\mathsf{C}}_{3}}\rtimes{{\mathsf{C}}_{4}}}
\left({\left\langle{a^2b}\right\rangle}\right) = 
{{{\mathsf{C}}_{3}}\rtimes{{\mathsf{C}}_{4}}}\\
\end{align*}

Los centralizadores de cada uno de los subgrupos en el subgrupo de estudio o ambiente son:

\begin{align*}
&{\mathsf{C}_{{{{\mathsf{C}}_{3}}\rtimes{{\mathsf{C}}_{4}}}}}
\left({\left\lbrace{1}\right\rbrace}\right) = 
{{{\mathsf{C}}_{3}}\rtimes{{\mathsf{C}}_{4}}}\\
&{\mathsf{C}_{{{{\mathsf{C}}_{3}}\rtimes{{\mathsf{C}}_{4}}}}}
\left({\left\langle{a^2}\right\rangle}\right) = 
{{{\mathsf{C}}_{3}}\rtimes{{\mathsf{C}}_{4}}}\\
&{\mathsf{C}_{{{{\mathsf{C}}_{3}}\rtimes{{\mathsf{C}}_{4}}}}}
\left({\left\langle{b}\right\rangle}\right) = 
{\left\langle{a^2b}\right\rangle}\\
&{\mathsf{C}_{{{{\mathsf{C}}_{3}}\rtimes{{\mathsf{C}}_{4}}}}}
\left({\left\langle{a}\right\rangle}\right) = 
{\left\langle{a}\right\rangle}\\
&{\mathsf{C}_{{{{\mathsf{C}}_{3}}\rtimes{{\mathsf{C}}_{4}}}}}
\left({\left\langle{ab}\right\rangle}\right) = 
{\left\langle{ab}\right\rangle}\\
&{\mathsf{C}_{{{{\mathsf{C}}_{3}}\rtimes{{\mathsf{C}}_{4}}}}}
\left({\left\langle{ab^2}\right\rangle}\right) = 
{\left\langle{ab^2}\right\rangle}\\
&{\mathsf{C}_{{{{\mathsf{C}}_{3}}\rtimes{{\mathsf{C}}_{4}}}}}
\left({\left\langle{a^2b}\right\rangle}\right) = 
{\left\langle{a^2b}\right\rangle}
\end{align*}

Los subgrupos derivados son:

\begin{align*}
&{\left\langle{a^2}\right\rangle}^{(1)} = {\left\lbrace{1}\right\rbrace}\\
&{\left\langle{b}\right\rangle}^{(1)} = {\left\lbrace{1}\right\rbrace}\\
&{\left\langle{a}\right\rangle}^{(1)} = {\left\lbrace{1}\right\rbrace}\\
&{\left\langle{ab}\right\rangle}^{(1)} = {\left\lbrace{1}\right\rbrace}\\
&{\left\langle{ab^2}\right\rangle}^{(1)} = {\left\lbrace{1}\right\rbrace}\\
&{\left\langle{a^2b}\right\rangle}^{(1)} = {\left\lbrace{1}\right\rbrace}\\
&{{{\mathsf{C}}_{3}}\rtimes{{\mathsf{C}}_{4}}}^{(1)} ={\left\langle{a^2}\right\rangle}\\
&{{{\mathsf{C}}_{3}}\rtimes{{\mathsf{C}}_{4}}}^{(2)} ={\left\lbrace{1}\right\rbrace}
\end{align*}

De aquí que el grupo que nos ocupa es resoluble.

Los subgrupos característicos que nos dicen si el grupo nilpotente:

\begin{align*}
&\mathsf{L}_{1}\left({\left\langle{a^2}\right\rangle}\right) 
= 
{\left\lbrace{1}\right\rbrace}\\
&\mathsf{L}_{1}\left({\left\langle{b}\right\rangle}\right) 
= 
{\left\lbrace{1}\right\rbrace}\\
&\mathsf{L}_{1}\left({\left\langle{a}\right\rangle}\right) 
= 
{\left\lbrace{1}\right\rbrace}\\
&\mathsf{L}_{1}\left({\left\langle{ab}\right\rangle}\right) 
= 
{\left\lbrace{1}\right\rbrace}\\
&\mathsf{L}_{1}\left({\left\langle{ab^2}\right\rangle}\right) 
= 
{\left\lbrace{1}\right\rbrace}\\
&\mathsf{L}_{1}\left({\left\langle{a^2b}\right\rangle}\right) 
= 
{\left\lbrace{1}\right\rbrace}\\
&{\mathsf{L}}_{1,2 ...}
\left(
{{{\mathsf{C}}_{3}}\rtimes{{\mathsf{C}}_{4}}}
\right) 
= 
{\left\langle{b}\right\rangle} \Rightarrow 
{{{\mathsf{C}}_{3}}\rtimes{{\mathsf{C}}_{4}}} 
\mbox{ NO ES NILPOTENTE}
\end{align*}

El subgrupo de Frattini de ${{\mathsf{C}}_{3}}\rtimes{{\mathsf{C}}_{4}}$ es 
$
\left\langle{ a^2 b}\right\rangle 
\cap 
\left\langle{a}\right\rangle 
\cap 
\left\langle{a b^2}\right\rangle
\cap
\left\langle{a b}\right\rangle
= 
{
\left\langle
{a^2}
\right\rangle
}$. 
Así: 
\[
{\Phi}_{{{{\mathsf{C}}_{3}}\rtimes{{\mathsf{C}}_{4}}}} = 
{\mathsf{Z}}{\left({{{\mathsf{C}}_{3}}\rtimes{{\mathsf{C}}_{4}}}\right)}=
{\left\langle{a^2}\right\rangle}
\]

Para el grupo de Fitting solo tenemos que coger los subgrupos nilpotentes normales, esto es, 
los subgrupos normales y y dar con el subgrupo generado por estos: ${\mathsf{F}}
{\left({{{\mathsf{C}}_{3}}\rtimes{{\mathsf{C}}_{4}}}\right)}=
{\left\langle{a^2 b}\right\rangle}$ ya que 
${\left\langle{a^2}\right\rangle}\subset{\left\langle{a^2 b}\right\rangle}$ 
y ${\left\langle{b}\right\rangle}\subset{\left\langle{a^2 b}\right\rangle}$, 
de forma que ${\left\langle{a^2 b}\right\rangle}={\left\langle{a^2 , b}\right\rangle}$.


Faltaría ver a qué equivale 
$\nicefrac{{{{\mathsf{C}}_{3}}\rtimes{{\mathsf{C}}_{4}}}}{\left\langle{a^2}\right\rangle}$
, de la misma forma que sería interesante ver 
$\nicefrac{{{{\mathsf{C}}_{3}}\rtimes{{\mathsf{C}}_{4}}}}{\left\langle{b}\right\rangle}$
y $\nicefrac{{{{\mathsf{C}}_{3}}\rtimes{{\mathsf{C}}_{4}}}}{\left\langle{a^2 b}\right\rangle}$.
Como el grupo del denominador es el centro cualquier conjugado constituye un clase de 
equivalencia que incluye $x,y,z$ no necesariamente distintas, talque, $xy=1 \wedge xz=a^2$:

\begin{align*}
&\overline{1}
&\triangleq&\quad
{1}\cdot{\left\langle{a^2}\right\rangle}
&\equiv&\quad
{\left\lbrace{1,a^2}\right\rbrace}\\
&\overline{a}
&\triangleq&\quad
{a}\cdot{\left\langle{a^2}\right\rangle}
&\equiv&\quad
{\left\lbrace{a,a^3}\right\rbrace}\\
&\overline{ab}
&\triangleq&\quad
{ab}\cdot{\left\langle{a^2}\right\rangle}
&\equiv&\quad
{\left\lbrace{ab,a^3b}\right\rbrace}\\
&\overline{ab^2}
&\triangleq&\quad
{ab^2}\cdot{\left\langle{a^2}\right\rangle}
&\equiv&\quad
{\left\lbrace{ab^2,a^3b^2}\right\rbrace}
\\
&\overline{b}
&\triangleq&\quad 
{b}\cdot{\left\langle{a^2}\right\rangle}
&\equiv&\quad
{\left\lbrace{b,a^2b}\right\rbrace}\\
&\overline{b^2}
&\triangleq&\quad 
{b^2}\cdot{\left\langle{a^2}\right\rangle}
&\equiv&\quad
{\left\lbrace{b^2,a^2b^2}\right\rbrace}
\end{align*}

En cuanto a la operación heredada, tenemos:

\begin{align*}
&\overline{1}\cdot\overline{x}=\overline{x}\cdot\overline{1}=\overline{x}\quad \mbox{conforma el grupo trivial}\\
&{\overline{b}}^{-1}={\overline{b^2}}\quad \mbox{conforman un grupo cíclico de orden 3}\\
&{\overline{a}}^{-1}={\overline{a}}\quad \mbox{conforman un grupo cíclico de orden 2}\\
&{\overline{ab}}^{-1}={\overline{ab}}\quad \mbox{conforman un grupo cíclico de orden 2}\\
&{\overline{ab^2}}^{-1}={\overline{ab^2}}\quad \mbox{conforman un grupo cíclico de orden 2}\\
&{\overline{b}}\cdot{\overline{a}}=\overline{ab^2}\\
&{\overline{a}}\cdot{\overline{b}}=\overline{ab}\\
&{\overline{b}}\cdot{\overline{a}}=\overline{ab^2}\neq\overline{ab}={\overline{a}}\cdot{\overline{b}}\\
&\nicefrac{{{\mathsf{C}}_{3}}\rtimes{{\mathsf{C}}_{4}}}{\left\langle a^2 \right\rangle} \cong{{\mathsf{Dih}}_{3}}\cong {\mathtt{Int}\left({{{\mathsf{C}}_{3}}\rtimes{{\mathsf{C}}_{4}}}\right)}
\end{align*}
 
En cuanto a todos los automorfismos sabemos que tienen que cumplir que $1 \mapsto 1 , a^2 \mapsto a^2$ y tenemos las posibilidades $a \mapsto x , \mathsf{o}(x)=4, ab \mapsto y, \mathsf{o}(y)=4$ y $b\mapsto b,b\mapsto b^2$ de forma que sean compatibles, que da un máximo de $12$ posibilidades. Como medida de control, todo automorfismo en el grupo total, cuando se restringe a $\langle a^2 b\rangle$ ha de ser también un automorfismo de $\mathtt{Aut}\left(\langle a^2 b\rangle\right)\cong\mathtt{Aut}\left({{\mathsf{C}}_{6}}\right)\cong {{\mathsf{C}}_{2}}$.

Por último también tenemos $\nicefrac{{{{\mathsf{C}}_{3}}\rtimes{{\mathsf{C}}_{4}}}}{\left\langle b \right\rangle} \cong {{\mathsf{C}}_{4}}$, y muy claramente , $\nicefrac{{{{\mathsf{C}}_{3}}\rtimes{{\mathsf{C}}_{4}}}}{\left\langle a^2 b \right\rangle} \cong {{\mathsf{C}}_{2}}$.

El diagrama de Hasse quedaría:


\begin{center}
    \begin{tikzpicture}
    % First, locate each of the nodes and name them
    \node (a-b) at (0,12)
    {$\mathsf{C}_{3}\rtimes\mathsf{C}_{4}\cong\left\langle{a,b\vert a^4=b^3=1,xa^2=a^2x,ab=b^2a}\right\rangle$};
    \node (a2b) at (3,10) 
    {$\mathsf{C}_{6}\cong \left\langle{a^2 b}\right\rangle$};
    \node (a) at (1,8) 
    {$\cong\left\langle{a}\right\rangle$};
    \node (ab) at (-1,8) 
    {$\cong\left\langle{ab}\right\rangle$};
    \node (ab2) at (-3,8) 
    {$\mathsf{C}_{4}\cong\left\langle{ab^2}\right\rangle$};
    \node (b) at (3,6) 
    {$\mathsf{C}_{3}\cong\left\langle{b}\right\rangle$};
    \node (a2) at (0,4) 
    {$\mathsf{C}_{2}\cong\left\langle{a^2}\right\rangle$};
    \node (1) at (0,0) 
    {${\mathsf{C}_{1}}\cong{\left\langle{1_{\mathsf{C}_{3}\rtimes\mathsf{C}_{4}}}\right\rangle}$};
   % Now draw the lines:
    \draw [black, thick] (a-b) -- (a2b);
    \draw [blue, thick] (a-b) -- (a);
    \draw [blue, thick] (a-b) -- (ab);
    \draw [blue, thick] (a-b) -- (ab2);
    \draw [black, thick] (a2b) -- (b);
    \draw [black, thick] (a2b) -- (a2);
    \draw [black, thick] (a) -- (a2);
    \draw [black, thick] (ab) -- (a2);
    \draw [black, thick] (ab2) -- (a2);
    \draw [black, thick] (b) -- (1);
    \draw [black, thick] (a2) -- (1);
    \end{tikzpicture}
\end{center}

El siguiente grupo no abeliano lo buscaremos también por los productos semidirectos de grupos de orden $4$ y de orden $3$. Se nos plantea el mismo problema que antes, ver si son posibles ${\left({{{\mathsf{C}}_{2}}\times{{\mathsf{C}}_{2}}}\right)}\rtimes{{\mathsf{C}}_{3}}$ y/o ${{\mathsf{C}}_{3}}\rtimes{\left({{{\mathsf{C}}_{2}}\times{{\mathsf{C}}_{2}}}\right)}$. Comenzamos por el último, esto es, si es posible construir un homomorfismo $\Psi : {{{\mathsf{C}}_{2}}\times{{\mathsf{C}}_{2}}} \longrightarrow {\mathtt{Aut}}{\left({{\mathsf{C}}_{3}}\right)}\cong{{\mathsf{C}}_{2}}$.

Veamos:
\begin{flushleft}
\begin{align*}
&{{\mathsf{C}}_{3}} &\triangleq \qquad&{\left\langle{c\mid c^3=1}\right\rangle}&&\\
&{{\mathsf{C}}_{2}} &\triangleq \qquad&{\left\langle{a,b\mid a^2=b^2=1, ab=ba}\right\rangle}&&\\
&{\mathtt{Aut}}{\left({{\mathsf{C}}_{3}}\right)} &\triangleq \qquad&\left\lbrace{\left\langle c\mapsto c \right\rangle,\left\langle c\mapsto c^2 \right\rangle}\right\rbrace&&
\end{align*}

\begin{align*}
&\Psi_0 : &{{{\mathsf{C}}_{2}}\times{{\mathsf{C}}_{2}}}\qquad &\longrightarrow &{\mathtt{Aut}}{\left({{\mathsf{C}}_{3}}\right)}\cong{{\mathsf{C}}_{2}}&\\
&&a&\mapsto\qquad & \left\langle c\mapsto c \right\rangle&\\
&&b&\mapsto\qquad & \left\langle c\mapsto c \right\rangle&\\
&&1&\mapsto\qquad & \left\langle c\mapsto c \right\rangle &\mbox{ en todo homomorfismo}\\
&&ab&\mapsto\qquad & \left\langle c\mapsto c \right\rangle &\mbox{ conclusión}\\
&\Psi_1 : &{{{\mathsf{C}}_{2}}\times{{\mathsf{C}}_{2}}}\qquad &\longrightarrow &{\mathtt{Aut}}{\left({{\mathsf{C}}_{3}}\right)}\cong{{\mathsf{C}}_{2}}&\\
&&a&\mapsto\qquad & \left\langle c\mapsto c^2 \right\rangle&\\
&&b&\mapsto\qquad & \left\langle c\mapsto c^2 \right\rangle&\\
&&1&\mapsto\qquad & \left\langle c\mapsto c \right\rangle &\mbox{ en todo homomorfismo}\\ 
&&ab&\mapsto\qquad & \left\langle c\mapsto c \right\rangle &\mbox{ conclusión}
\end{align*}
\end{flushleft}

Tomando $\Phi_1$ como nuestro homomorfismo, puesto que $\Phi_0$ no es mas que la identidad, que nos daría un producto directo y por lo tanto abeliano. La acción quedaría así:

\begin{flushright}
\begin{align*}
&\Psi_1 : &{{{\mathsf{C}}_{2}}\times{{\mathsf{C}}_{2}}}\times{{\mathsf{C}}_{3}} &\longrightarrow &{{\mathsf{C}}_{3}}\\
&&a^ib^jc^k&\mapsto\qquad & c^l
\end{align*}


\begin{align*}
&l=
\begin{cases}
-k \Rightarrow (i\neq j)\\
k \;\Rightarrow (i=j)
\end{cases}
\\
& i,j \in \left\lbrace 0,1 \right\rbrace\\ 
& k,l \in \left\lbrace 0,1,2 \right\rbrace
\end{align*}
\end{flushright}

	Aplicando esta acción al producto directo sale la siguiente tabla de multiplicar:


\begin{align*}
\begin{array}{|c|c|c|c|c|c|c|c|c|c|c|c|}
\hline
1&g&g^2&b&bg&bg^2&a&ag&ag^2&ab&abg&abg^2\\
\hline
g&g^2&1&bg^2&b&bg&ag^2&a&ag&abg&abg^2&ab\\
\hline
g^2&1&g&bg&bg^2&b&ag&ag^2&a&abg^2&ab&abg\\
\hline
b&bg&bg^2&1&g&g^2&ab&abg&abg^2&a&ag&ag^2\\
\hline
bg&bg^2&b&g^2&1&g&abg^2&ab&abg&ag&ag^2&a\\
\hline
bg^2&b&bg&g&g^2&1&abg&abg^2&ab&ag^2&a&ag\\
\hline
a&ag&ag^2&ab&abg&abg^2&1&g&g^2&b&bg&bg^2\\
\hline
ag&ag^2&a&abg^2&ab&abg&g^2&1&g&bg&bg^2&b\\
\hline
ag^2&a&ag&abg&abg^2&ab&g&g^2&1&bg^2&b&bg\\
\hline
ab&abg&abg^2&a&ag&ag^2&b&bg&bg^2&1&g&g^2\\
\hline
abg&abg^2&ab&ag^2&a&ag&bg^2&b&bg&g&g^2&1\\
\hline
abg^2&ab&abg&ag&ag^2&a&bg&bg^2&b&g^2&1&g\\
\hline
\end{array}
\end{align*}

Primero vemos que el orden de los elementos es:

\begin{align*}
&\mathsf{o}{\left({1}\right)} = 1\quad &&\mathsf{o}{\left({g}\right)} = 3\quad &&\mathsf{o}{\left({g^2}\right)} = 3\\
&\mathsf{o}{\left({b}\right)} = 2\quad &&\mathsf{o}{\left({bg}\right)} = 2\quad &&\mathsf{o}{\left({bg^2}\right)} = 2\\
&\mathsf{o}{\left({a}\right)} = 2\quad &&\mathsf{o}{\left({ag}\right)} = 2\quad &&\mathsf{o}{\left({ag^2}\right)} = 2\\
&\mathsf{o}{\left({ab}\right)} = 2\quad &&\mathsf{o}{\left({abg}\right)} = 6\quad &&\mathsf{o}{\left({abg^2}\right)} = 6
\end{align*}

De aquí ya podemos decir que subgrupos cíclicos tiene nuestro grupo:

\begin{align*}
&{\mathcal{L}_{1}}{\left({G}\right)}\triangleq{\left\lbrace{\langle x \rangle\vert x\in G}\right\rbrace}\\
&\\ 
&{\mathcal{L}_{1}}{\left({{\mathsf{C}_{3}}\rtimes{\mathsf{C}_{2}\times\mathsf{C}_{2}}}\right)} = {\left\lbrace\right.}
{\left\lbrace{1}\right\rbrace},{\left\lbrace{1,b}\right\rbrace},\\
&\qquad\qquad\qquad\qquad{\left\lbrace{1,bg}\right\rbrace},
{\left\lbrace{1,bg^2}\right\rbrace},\\
&\qquad\qquad\qquad\qquad{\left\lbrace{1,a}\right\rbrace},
{\left\lbrace{1,ag}\right\rbrace},
{\left\lbrace{1,ag^2}\right\rbrace},
{\left\lbrace{1,ag}\right\rbrace},\\
&\qquad\qquad\qquad\qquad{\left\lbrace{1,g,g^2}\right\rbrace},
{\left\lbrace{1,g,g^2,ab,abg,abg^2}\right\rbrace}{\left.\right\rbrace}
\end{align*}

Los subgrupos generados estrictamente por $2$ elementos son:

\begin{align*}
&{\mathcal{L}_{2}}{\left({G}\right)}\triangleq
{\left\lbrace{
\langle x,y \rangle
\vert 
(x,y\in G) \wedge (\nexists z \in G \;\langle x,y \rangle = \langle z\rangle)
}\right\rbrace}\\
&\\
&{\mathcal{L}_{2}}{\left({{\mathsf{C}_3}\rtimes{({\mathsf{C}_{2}}\times{\mathsf{C}_{2}})}}\right)}=\left\lbrace\right.
{\left\lbrace{1,b,a,ab}\right\rbrace},
{\left\lbrace{1,bg,ag,ab}\right\rbrace},
{\left\lbrace{1,bg^2,ag^2,ab}\right\rbrace},\\
&\qquad\qquad\qquad\qquad\qquad\;\,{\left\lbrace{1,g,g^2,b,bg,bg^2,a,ag,ag^2,ab,abg,abg^2}\right\rbrace}
\left.\right\rbrace
\end{align*}

Con esto queda claro que no hay subgrupos propiamente generados por $3$ elementos de forma minimal, pues con $2$ elementos se genera todo el grupo.

El centro del grupo es: 
\[
\mathsf{Z}\left({{\mathsf{C}_{3}}\rtimes({\mathsf{C}_{2}}\times{\mathsf{C}_{2}})}\right) = {\left\lbrace{1,ab}\right\rbrace}
\]


Los subgrupos que son " absolutamente no normales ", esto es, aquellos que coinciden con su normalizador:


\begin{align*}
&
\mathsf{N}_{{\mathsf{C}_{3}}\rtimes({\mathsf{C}_{2}}\times{\mathsf{C}_{2}})}
\left(
{\left\lbrace
{1,b,a,ab}
\right\rbrace}
\right) = 
{\left\lbrace
{1,b,a,ab}
\right\rbrace}\\
&
\mathsf{N}_{{\mathsf{C}_{3}}\rtimes({\mathsf{C}_{2}}\times{\mathsf{C}_{2}})}
\left(
{\left\lbrace
{1,bg,ag,ab}
\right\rbrace}
\right) = 
{\left\lbrace
{1,bg,ag,ab}
\right\rbrace}\\
&
\mathsf{N}_{{\mathsf{C}_{3}}\rtimes({\mathsf{C}_{2}}\times{\mathsf{C}_{2}})}
\left(
{\left\lbrace
{1,bg^2,ag^2,ab}
\right\rbrace}
\right) = 
{\left\lbrace
{1,bg^2,ag^2,ab}
\right\rbrace}
\end{align*}

A continuación, los que tienen un normalizador estrictamente mayor que ellos mismos, pero sin ser normales, son:

\begin{align*}
&\mathsf{N}_{{\mathsf{C}_{3}}\rtimes({\mathsf{C}_{2}}\times{\mathsf{C}_{2}})}
\left({\left\lbrace{1,b}\right\rbrace}\right) =
\mathsf{N}_{{\mathsf{C}_{3}}\rtimes({\mathsf{C}_{2}}\times{\mathsf{C}_{2}})}
\left({\left\lbrace{1,a}\right\rbrace}\right) = 
\left\lbrace{1,b,a,ab}\right\rbrace
\\
&\mathsf{N}_{{\mathsf{C}_{3}}\rtimes({\mathsf{C}_{2}}\times{\mathsf{C}_{2}})}
\left({\left\lbrace{1,bg}\right\rbrace}\right) = 
\mathsf{N}_{{\mathsf{C}_{3}}\rtimes({\mathsf{C}_{2}}\times{\mathsf{C}_{2}})}
\left({\left\lbrace{1,ag}\right\rbrace}\right) = 
{\left\lbrace{1,bg,ag,ab}\right\rbrace}
\\
&\mathsf{N}_{{\mathsf{C}_{3}}\rtimes({\mathsf{C}_{2}}\times{\mathsf{C}_{2}})}
\left({\left\lbrace{1,bg^2}\right\rbrace}\right) = 
\mathsf{N}_{{\mathsf{C}_{3}}\rtimes({\mathsf{C}_{2}}\times{\mathsf{C}_{2}})}
\left({\left\lbrace{1,ag^2}\right\rbrace}\right) = 
{\left\lbrace{1,bg^2,ag^2,ab}\right\rbrace}
\end{align*}

Los subgrupos normales quedan:

\begin{align*}
&{\mathcal{L}}_{{{\mathsf{C}}_{3}}\rtimes
{\left({
{{\mathsf{C}}_{2}}\times{{\mathsf{C}}_{2}}
}\right)}}=
{\left\lbrace\right.}
{\left\lbrace{1}\right\rbrace},
{\left\lbrace{1,ab}\right\rbrace},
{\left\lbrace{1,g,g^2}\right\rbrace},
{\left\lbrace{1,g,g^2,b,bg,bg^2}\right\rbrace},\\
&{\left\lbrace{1,g,g^2,a,ag,ag^2}\right\rbrace},
{\left\lbrace{1,g,g^2,ab,abg,abg^2}\right\rbrace},
{{\mathsf{C}}_{3}}
\rtimes
{\left({
{{\mathsf{C}}_{2}}\times{{\mathsf{C}}_{2}}
}\right)}
{\left.\right\rbrace}
\end{align*}

De estos hay $2$ subgrupos no abelianos de orden $6$ que no son nilpotentes, y de hecho son isomorfos a $\mathsf{Dih}_{3}$, a saber son ${\left\lbrace{1,g,g^2,b,bg,bg^2}\right\rbrace}$y ${\left\lbrace{1,g,g^2,a,ag,ag^2}\right\rbrace}$. Por lo tanto el retículo de subgrupos nilpotentes normales es:

\begin{align*}
&{\mathcal{L}}_{{{\mathsf{C}}_{3}}\rtimes
{\left({
{{\mathsf{C}}_{2}}\times{{\mathsf{C}}_{2}}
}\right)}}^{\mathcal{NIL}}=
\left\lbrace
{\left\lbrace{1}\right\rbrace},
{\left\lbrace{1,ab}\right\rbrace},
{\left\lbrace{1,g,g^2}\right\rbrace},
{\left\lbrace{1,g,g^2,ab,abg,abg^2}\right\rbrace}
\right\rbrace\\
&\left\langle{\bigcup{\mathcal{L}}_{{{\mathsf{C}}_{3}}\rtimes
{\left({
{{\mathsf{C}}_{2}}\times{{\mathsf{C}}_{2}}
}\right)}}^{\mathcal{NIL}}}\right\rangle = {\left\lbrace{1,g,g^2,ab,abg,abg^2}\right\rbrace} = \left\langle abg \right\rangle
\end{align*}

Buscamos ahora los centralizadores de cada subgrupo en el grupo ambiente, quitando el trivial, el impropio y el centro que ya se ha hallado.
Primero buscamos los subgrupos propios abelianos (lo son todos los propios) pero que su centralizador es el mismo subgrupo:

\begin{align*}
&\mathsf{C}_{{\mathsf{C}_{3}}\rtimes({\mathsf{C}_{2}}\times{\mathsf{C}_{2}})}\left({\left\lbrace{1,b,a,ab}\right\rbrace}\right) = {\left\lbrace{1,b,a,ab}\right\rbrace}\\
&\mathsf{C}_{{\mathsf{C}_{3}}\rtimes({\mathsf{C}_{2}}\times{\mathsf{C}_{2}})}\left({\left\lbrace{1,bg,ag,ab}\right\rbrace}\right) = {\left\lbrace{1,bg,ag,ab}\right\rbrace}\\
&\mathsf{C}_{{\mathsf{C}_{3}}\rtimes({\mathsf{C}_{2}}\times{\mathsf{C}_{2}})}\left({\left\lbrace{1,bg^2,ag^2,ab}\right\rbrace}\right) = {\left\lbrace{1,bg^2,ag^2,ab}\right\rbrace}\\
&\mathsf{C}_{{\mathsf{C}_{3}}\rtimes({\mathsf{C}_{2}}\times{\mathsf{C}_{2}})}\left({\left\lbrace{1,g,g^2,ab,abg,abg^2}\right\rbrace}\right) = {\left\lbrace{1,g,g^2,ab,abg,abg^2}\right\rbrace}
\end{align*}

A continuación aquellos subgrupos que tienen centralizador mayor que ellos mismos:

\begin{align*}
&\mathsf{C}_{{\mathsf{C}_{3}}\rtimes
({\mathsf{C}_{2}}\times{\mathsf{C}_{2}})}\left({\left\lbrace{1,b}\right\rbrace}\right) =
\mathsf{C}_{{\mathsf{C}_{3}}\rtimes
({\mathsf{C}_{2}}\times{\mathsf{C}_{2}})}\left({\left\lbrace{1,a}\right\rbrace}\right) = 
{\left\lbrace{1,b,a,ab}\right\rbrace}\\
&\mathsf{C}_{{\mathsf{C}_{3}}\rtimes
({\mathsf{C}_{2}}\times{\mathsf{C}_{2}})}\left({\left\lbrace{1,bg}\right\rbrace}\right) =
\mathsf{C}_{{\mathsf{C}_{3}}
\rtimes({\mathsf{C}_{2}}\times{\mathsf{C}_{2}})}\left({\left\lbrace{1,ag}\right\rbrace}\right) = 
{\left\lbrace{1,bg,ag,ab}\right\rbrace}\\
&\mathsf{C}_{{\mathsf{C}_{3}}
\rtimes
({\mathsf{C}_{2}}\times{\mathsf{C}_{2}})}\left({\left\lbrace{1,bg^2}\right\rbrace}\right) = 
\mathsf{C}_{{\mathsf{C}_{3}}
\rtimes
({\mathsf{C}_{2}}\times{\mathsf{C}_{2}})}\left({\left\lbrace{1,ag^2}\right\rbrace}\right) = {\left\lbrace{1,bg^2,ag^2,ab}\right\rbrace}\\
&\mathsf{C}_{{\mathsf{C}_{3}}\rtimes({\mathsf{C}_{2}}\times{\mathsf{C}_{2}})}\left({\left\lbrace{1,g,g^2,abg}\right\rbrace}\right) = {\left\lbrace{1,g,g^2,ab,abg,abg^2}\right\rbrace}
\end{align*}


Los subgrupos derivados vienen:

\begin{align*}
&\mathsf{D^{(1)}}{\left({\mathsf{C}_{3}}\rtimes({\mathsf{C}_{2}}\times{\mathsf{C}_{2}})\right)} = {\left\lbrace{1,g,g^2}\right\rbrace}\\
&\mathsf{D^{(2)}}{\left({\mathsf{C}_{3}}\rtimes({\mathsf{C}_{2}}\times{\mathsf{C}_{2}})\right)} = {\left\lbrace{1}\right\rbrace}
\end{align*}

De dónde vemos que el grupo ${\mathsf{C}_{3}}\rtimes({\mathsf{C}_{2}}\times{\mathsf{C}_{2}})$ \textsc{ es resoluble}.

Los subgrupos L-derivados del grupo máximo (en el retículo):

\begin{align*}
\mathsf{L_{(1)}}{\left({\mathsf{C}_{3}}\rtimes({\mathsf{C}_{2}}\times{\mathsf{C}_{2}})
\right)} = \mathsf{L_{(2)}}{\left({\mathsf{C}_{3}}\rtimes({\mathsf{C}_{2}}\times{\mathsf{C}_{2}})
\right)} = {\left\lbrace{1,g,g^2}\right\rbrace}
\end{align*}

De dónde vemos que el grupo 
${\mathsf{C}_{3}}\rtimes({\mathsf{C}_{2}}\times{\mathsf{C}_{2}})$ \textsc{no es nilpotente}.

Para aclararnos un poco entre tanto elemento y subgrupo veamos como se pueden generar los 
subgrupos:

\begin{align*}
&{\left\langle{b}\right\rangle} = {\left\lbrace{1,b}\right\rbrace}\\
&{\left\langle{bg}\right\rangle} = {\left\lbrace{1,bg}\right\rbrace}\\
&{\left\langle{bg^2}\right\rangle} = {\left\lbrace{1,bg^2}\right\rbrace}\\
&{\left\langle{a}\right\rangle} = {\left\lbrace{1,a}\right\rbrace}\\
&{\left\langle{ag}\right\rangle} = {\left\lbrace{1,ag}\right\rbrace}\\
&{\left\langle{ag^2}\right\rangle} = {\left\lbrace{1,ag^2}\right\rbrace}\\
&{\left\langle{ab}\right\rangle} = {\left\lbrace{1,ab}\right\rbrace}\\
&{\left\langle{g,b}\right\rangle} = {\left\langle{g^2,b}\right\rangle} = {\left\langle{g,bg}\right\rangle} = {\left\langle{g,bg^2}\right\rangle} = {\left\lbrace{1,g,g^2,b,bg,bg^2}\right\rbrace}\\
&{\left\langle{g,a}\right\rangle} = {\left\langle{g^2,a}\right\rangle} ={\left\langle{g,ag}\right\rangle} = {\left\langle{g,ag^2}\right\rangle} = {\left\lbrace{1,g,g^2,a,ag,ag^2}\right\rbrace}\\
&{\left\langle{g}\right\rangle} = {\left\lbrace{1,g,g^2}\right\rbrace}\\
&{\left\langle{abg}\right\rangle} =
{\left\langle{g,ab}\right\rangle} = {\left\langle{g^2,ab}\right\rangle} =  
{\left\lbrace{1,g,g^2,ab,abg,abg^2}\right\rbrace}\\
&{\left\langle{b,a}\right\rangle} = {\left\langle{b,ab}\right\rangle} = {\left\lbrace{1,b,a,ab}\right\rbrace}\\
&{\left\langle{bg,ag}\right\rangle} = {\left\langle{bg,ab}\right\rangle} = {\left\lbrace{1,bg,ag,ab}\right\rbrace}\\
&{\left\langle{bg^2,ag^2}\right\rangle} = {\left\langle{bg^2,ab}\right\rangle} = {\left\lbrace{1,bg^2,ag^2,ab}\right\rbrace}\\
&{\left\langle{g,b,a}\right\rangle} = {\left\lbrace{1,g,g^2,b,bg,bg^2,a,ag,ag^2,ab,abg,abg^2}\right\rbrace}
\end{align*}

El subgrupo de Frattini, la intersección de los subgrupos maximales, es el subgrupo trivial, ${\Phi}_{\mathsf{C}_{3}}\rtimes{\left({\mathsf{C}_{2}}\times{\mathsf{C}_{2}}\right)}= \lbrace 1\rbrace$. El subgrupo de Fitting, la suma de todos los subgrupos nilpotentes (todos excepto el propio) normales, coincidirá con el ya hallado anteriormente, $\mathsf{F}\left(\mathsf{C}_{3}\rtimes\left({\mathsf{C}_{2}}\times{\mathsf{C}_{2}}\right)\right)=\langle abg \rangle$.

Solo nos queda ver como se factoriza el grupo, primero a través del centro:
 
\begin{align*}
&
\nicefrac
{\mathsf{C}_{3}\rtimes\left({\mathsf{C}_{2}}\times{\mathsf{C}_{2}}\right)}
{\left\langle ab \right\rangle}
=
\left\lbrace\right.
1{\left\langle{ab}\right\rangle},
g{\left\langle{ab}\right\rangle},
g^2{\left\langle{ab}\right\rangle},
b{\left\langle{ab}\right\rangle}=
a{\left\langle{ab}\right\rangle},
bg{\left\langle{ab}\right\rangle}
=
ag{\left\langle{ab}\right\rangle},\\
&
bg^2{\left\langle{ab}\right\rangle}
=
ag^2{\left\langle{ab}\right\rangle}
\left.\right\rbrace
\end{align*}

Que tomando representantes:

\begin{align*}
&
\nicefrac
{\mathsf{C}_{3}\rtimes\left({\mathsf{C}_{2}}\times{\mathsf{C}_{2}}\right)}
{\left\langle ab \right\rangle}
=
\left\lbrace
\overline{1},
\overline{g},
\overline{g^2},
\overline{a},
\overline{ag},
\overline{ag^2}
\right\rbrace
\end{align*}

Que tiene la estructura del grupo $\mathsf{Dih}_{3}$ y el completo tiene la apariencia de $\mathsf{Dih}_{3}\times\mathsf{C}_{2}$. 

Para empezar, $ab$ conmuta con todo el grupo, y $\langle a,g \rangle\cong\mathsf{Dih}_{3}$ es también grupo normal. De otro lado $\langle ab\rangle \cap \langle a,g \rangle = \lbrace 1 \rbrace$. Por lo tanto: 
\[
{\mathsf{C}_{3}\rtimes\left({\mathsf{C}_{2}\times\mathsf{C}_{2}}\right)}\cong{\mathsf{Dih}_{3}\times\mathsf{C}_{2}}
\]

Si llamamos a $\mathtt{f}\triangleq abg$, $\mathtt{g}\triangleq a$ y $\mathtt{a} \triangleq ab$ explicitamos el isomorfismo antes dicho. 

Observamos con facilidad que ${\mathsf{C}_{3}}\rtimes{\left({{\mathsf{C}_{2}}\times{\mathsf{C}_{2}}}\right)}\cong\mathsf{Dih}_{6}$. Veamos esto último. Sustituimos $abg$ por $\mathfrak{f}\triangleq abg$, siendo $(abg)^2 = g^2$ de dónde  $\mathfrak{f}^{2}\triangleq g^2$, ${(abg)}^{3}=ab$, obtenemos $\mathfrak{f}^{3}\triangleq ab$, $(abg)^{4}= g$ que nos da $\mathfrak{f}^{4}\triangleq g$ y por último $(abg)^{5}= abg^2$ que se traslada a $\mathfrak{f}^{5}\triangleq abg^2$.  Podemos escoger $b$ como elemento no normal y de orden $2$ del grupo dihédrico, escogeremos $\mathfrak{g}\triangleq b$. Como resultado $\mathfrak{gf}=a$, y tenemos solo $2$ generadores como habíamos deducido al generar con la herramienta el grupo total. Via las sustituciones hechas, vemos que:
 
\[{\mathsf{C}_{3}}\rtimes{\left({{\mathsf{C}_{2}}\times{\mathsf{C}_{2}}}\right)}\cong{\mathsf{C}_{6}}\rtimes{\mathsf{C}_{2}}\cong\mathsf{Dih}_{6}\cong{\mathsf{Dih}_{3}}\rtimes{\mathsf{C}_{2}}\] 

Vemos como factoriza a través de $\langle abg\rangle$:

\begin{align*}
\overline{a}\triangleq{\left\lbrace{a,b,ag,ag^2,bg,bg^2}\right\rbrace},
\overline{1}\triangleq{\left\lbrace{1,g^2,g,ab,abg,abg^2}\right\rbrace}
\end{align*}

Que no es otra cosa que el grupo cíclico de orden $2$. Por eso no tiene sentido hacer el cociente por los otros subgrupos de orden $6$, puesto que partírán el grupo en dos partes, la unidad y un elemento de orden $2$.

El diagrama de Hasse es notablemente más complicado que  los anteriores:


\begin{center}
    \begin{tikzpicture}
    % First, locate each of the nodes and name them
    \node (a-b-g) at (-8,12)
    {$\mathsf{Dih}_{6}
    \cong
    {\left\langle{a,b,g\vert a^2=b^2=g^3=1,xab=abx,ag=g^2a,bg=g^b}\right\rangle}$};
    \node (b-g) at (-14,9) 
    {$\left\langle{b,g\vert b^2=g^3=1,bg=g^2b}\right\rangle$};
    \node (a-g) at (-8,9) 
    {$\left\langle{a,g\vert a^2=g^3=1,ag=g^2a}\right\rangle$};
    \node (abg) at (-2,9) 
    {$\cong\left\langle{abg\vert \left(ab\right)^6=1}\right\rangle$};
    \node (a-b) at (-14,6) 
    {$\left\langle{a,b\vert a^2=b^2=1}\right\rangle$};
    \node (ag-bg) at (-8,6) 
	{$\left\langle{ag,bg\vert (ag)^2=(bg)^2=1}\right\rangle$};
    \node (ag2-bg2) at (-2,6) 
    {$\left\langle{ag^2,bg^2\vert (ag^2)^2=(bg^2)^2=1}\right\rangle$};
    \node (g) at (-8,3) 
	{$\left\langle{g\vert g^3=1}\right\rangle$};    
    \node (ab) at (-16,0) 
    {$\left\langle{ab}\right\rangle$};
    \node (b) at (-13,0) 
    {$\left\langle{b}\right\rangle$};
    \node (a) at (-10,0) 
    {$\left\langle{a}\right\rangle$};
    \node (bg) at (-7,0) 
	{$\left\langle{bg}\right\rangle$};
    \node (ag) at (-4,0) 
    {$\left\langle{ag}\right\rangle$};
    \node (bg2) at (-1,0) 
    {$\left\langle{bg^2}\right\rangle$};
    \node (ag2) at (2,0) 
    {$\left\langle{ag^2}\right\rangle$};
	\node (1) at (-8,-3) 
    {$\left\langle{1}\right\rangle$};
   % Now draw the lines:
    \draw [black, thick] (a-b-g) -- (abg);
    \draw [black, thick] (a-b-g) -- (a-g);
    \draw [black, thick] (a-b-g) -- (b-g);
    \draw [black, thick] (a-b-g) -- (a-b);
    \draw [black, thick] (a-b-g) -- (ag-bg);
    \draw [black, thick] (a-b-g) -- (ag2-bg2);
    \draw [black, thick] (abg) -- (g);
    \draw [black, thick] (abg) -- (ab);
    \draw [black, thick] (a-g) -- (a);
    \draw [black, thick] (a-g) -- (ag);
    \draw [black, thick] (a-g) -- (ag2);
    \draw [black, thick] (a-g) -- (g);
    \draw [black, thick] (b-g) -- (b);
    \draw [black, thick] (b-g) -- (bg);
    \draw [black, thick] (b-g) -- (bg2);
    \draw [black, thick] (b-g) -- (g);
    \draw [black, thick] (a-b) -- (a);
    \draw [black, thick] (a-b) -- (ab);
    \draw [black, thick] (a-b) -- (b);
    \draw [black, thick] (ag-bg) -- (ag);
    \draw [black, thick] (ag-bg) -- (bg);
    \draw [black, thick] (ag2-bg2) -- (ag2);
    \draw [black, thick] (ag2-bg2) -- (bg2);
    \draw [black, thick] (g) -- (1);
    \draw [black, thick] (a) -- (1);
    \draw [black, thick] (b) -- (1);
    \draw [black, thick] (ab) -- (1);
    \draw [black, thick] (ag) -- (1);
    \draw [black, thick] (bg) -- (1);
    \draw [black, thick] (ag2) -- (1);
    \draw [black, thick] (bg2) -- (1);
    \end{tikzpicture}
\end{center}




\end{document}
